% Problems, Chapter 12: Further Topics in Analysis
% Arfken, Weber & Harris, Mathematical Methods for Physicists, 7th ed.
% Requires: amsmath, amssymb, tikz

\section*{Problems}

\Needspace{8\baselineskip}
\subsection*{12.1 Orthogonal Polynomials}

\Needspace{5\baselineskip}
\par\bigskip
\noindent\textbf{12.1.1}\quad Starting from the Rodrigues formula in
Table 12.1 for the Hermite polynomials $H_{n}$, derive the generating
function for the $H_{n}$ given in that table.
\solution
The convention in Table 12.1 is
\[
H_n(x)=(-1)^ne^{x^2}\frac{d^n}{dx^n}e^{-x^2}.
\]
Taylor's theorem applied to the entire function $f(x)=e^{-x^2}$ gives
\[
f(x-t)=\sum_{n=0}^{\infty}\frac{(-t)^n}{n!}f^{(n)}(x).
\]
Multiply by $e^{x^2}$ and insert the Rodrigues expression term by term:
\[
\sum_{n=0}^{\infty}H_n(x)\frac{t^n}{n!}
=e^{x^2}e^{-(x-t)^2}
=\boxed{e^{2xt-t^2}}.
\]
The Taylor series converges for every finite complex $t$, so the generating
function is an identity, not only a formal expansion. As a normalization
check, its first three coefficients are $1,2x,2x^2-1$, corresponding to
$H_0=1$, $H_1=2x$, $H_2=4x^2-2$ after multiplying by $n!$.

\Needspace{5\baselineskip}
\par\bigskip
\noindent\textbf{12.1.2}\quad
\begin{enumerate}
\item[(a)] Starting from the Laguerre ODE,
  \[
  xy''+(1-x)y'+\lambda y=0,
  \]
  obtain the Rodrigues formula for its polynomial solutions $L_{n}(x)$.
\item[(b)] From the Rodrigues formula, scaled as in Table 12.1, derive the
  generating function for the $L_{n}(x)$ given in that table.
\end{enumerate}
\solution
(a) If $y$ is a polynomial of degree $n$ with leading coefficient $a_n\ne0$,
the coefficient of $x^n$ in the ODE is $(\lambda-n)a_n$. Thus
$\lambda=n$. The weight needed to put the equation in self-adjoint form
satisfies
\[
(xw)'=(1-x)w,\quad xw'+w=w-xw,\quad w'=-w,
\]
so choose $w=e^{-x}$. The Rodrigues construction with $p=x$ gives
\[
\boxed{L_n(x)=\frac{e^x}{n!}\frac{d^n}{dx^n}(x^ne^{-x})}.
\]
We can verify this construction directly. Put $F=x^ne^{-x}$, so
$xF'=(n-x)F$. Differentiate $n+1$ times. Since both $x$ and $n-x$ are
linear, the product rule leaves
\[
xF^{(n+2)}+(n+1)F^{(n+1)}
=(n-x)F^{(n+1)}-(n+1)F^{(n)}.
\]
With $G=F^{(n)}$ this becomes $xG''+(x+1)G'+(n+1)G=0$.
Now $G=n!e^{-x}L_n$, so
$G'=n!e^{-x}(L_n'-L_n)$ and
$G''=n!e^{-x}(L_n''-2L_n'+L_n)$.
Substitution cancels the common factor and gives
$xL_n''+(1-x)L_n'+nL_n=0$, exactly the required ODE.

Leibniz's rule also shows that the result is a polynomial:
\[
L_n(x)=\sum_{k=0}^n\binom nk\frac{(-x)^k}{k!}.
\]
Indeed, differentiating $x^n$ $n-k$ times gives $n!x^k/k!$, while the
remaining $k$ derivatives of $e^{-x}$ give $(-1)^ke^{-x}$.
This fixes the scale $L_n(0)=1$ and gives leading coefficient $(-1)^n/n!$.

(b) Sum that explicit polynomial formula:
\begin{align*}
\sum_{n=0}^{\infty}L_n(x)t^n
&=\sum_{k=0}^{\infty}\frac{(-x)^k}{k!}
       \sum_{n=k}^{\infty}\binom nk t^n\\
&=\sum_{k=0}^{\infty}\frac{(-x)^k}{k!}
       \frac{t^k}{(1-t)^{k+1}}\\
&=\boxed{\frac1{1-t}\exp\left(-\frac{xt}{1-t}\right)},\qquad |t|<1.
\end{align*}
The inner sum follows by differentiating $(1-t)^{-1}$ $k$ times and
multiplying by $t^k/k!$. Interchanging the sums is justified because the
corresponding absolute sum is at most
$(1-|t|)^{-1}\exp[|x||t|/(1-|t|)]<\infty$.

\Needspace{5\baselineskip}
\par\bigskip
\noindent\textbf{12.1.3}\quad Carry out in detail the steps needed to confirm
that the $(n+1)$-fold differentiation of Eq.~(12.10) leads to Eq.~(12.12).
\solution
Set $F=wp^n$ and $A=(n-1)p'+q$. Equation (12.10) is $pF'=AF$.
Because $p$ has degree at most two, $p^{(j)}=0$ for $j\ge3$; because
$A$ is at most linear, $A^{(j)}=0$ for $j\ge2$. Leibniz's formula
therefore gives, with every surviving term displayed,
\begin{align*}
D^{n+1}(pF')
 &=pF^{(n+2)}+(n+1)p'F^{(n+1)}
      +\frac{n(n+1)}2p''F^{(n)},\\
D^{n+1}(AF)
 &=AF^{(n+1)}+(n+1)A'F^{(n)}.
\end{align*}
The last term in the first line has coefficient zero if $n=0$.
Bring the right side to the left and divide by nonzero $w$.
The coefficient of $F^{(n+1)}/w$ is
\[
(n+1)p'-[(n-1)p'+q]=2p'-q.
\]
The coefficient of $F^{(n)}/w=y_n$ is
\begin{align*}
\frac{n(n+1)}2p''-(n+1)[(n-1)p''+q']
&=\frac{(n+1)(2-n)}2p''-(n+1)q'\\
&=-\left[\frac{n^2-n-2}{2}p''+(n+1)q'\right].
\end{align*}
Hence
\[
\boxed{\frac p wF^{(n+2)}+\frac{2p'-q}{w}F^{(n+1)}
-\left[\frac{n^2-n-2}{2}p''+(n+1)q'\right]y_n=0},
\]
which is Eq. (12.12).

\Needspace{5\baselineskip}
\par\bigskip
\noindent\textbf{12.1.4}\quad Confirm the algebraic steps that convert
Eq.~(12.12) into Eq.~(12.16).
\solution
Write $G=D^n(wp^n)=wy_n$ and
\[
S=\frac{w'}w=\frac{q-p'}p.
\]
Two product differentiations give
\[
\frac{G'}w=y_n'+Sy_n,\qquad
\frac{G''}w=y_n''+2Sy_n'+(S'+S^2)y_n.
\]
Put $A=(n^2-n-2)p''/2+(n+1)q'$. Substitution into Eq. (12.12) gives
\[
p y_n''+[2pS+2p'-q]y_n'
+[pS'+pS^2+(2p'-q)S-A]y_n=0.
\]
The first-derivative coefficient simplifies to
$2(q-p')+2p'-q=q$. In the last coefficient,
\[
pS^2+(2p'-q)S=S[(q-p')+2p'-q]=p'S,
\]
so
\[
pS'+p'S-A=(pS)'-A=q'-p''-A
=-\frac{n(n-1)}2p''-nq'.
\]
Thus
\[
\boxed{p y_n''+q y_n'-\left[\frac{n(n-1)}2p''+nq'\right]y_n=0},
\]
exactly Eq. (12.16), with eigenvalue
$\lambda_n=-n(n-1)\alpha-n\mu$ when $p''=2\alpha$, $q'=\mu$.

\erratum{The source's intermediate Eqs. (12.13)--(12.15), on textbook
p.~553, contain algebraic errors. Equation (12.16) is correct. The
consistent intermediate identities are the following:}
\begin{align*}
\erratum{\frac p wG''}
&\erratum{=p y_n''+\frac{2(q-p')}{w}G'
 -\left[p''-q'+\frac{q(q-p')}{p}\right]y_n,}\\
\erratum{0}
&\erratum{=p y_n''+\frac q wG'
 -\left[\frac{n(n-1)}2p''+nq'+\frac{q(q-p')}{p}\right]y_n,}\\
\erratum{\frac q wG'}
&\erratum{=q y_n'+\frac{q(q-p')}{p}y_n.}
\end{align*}
For the first identity, eliminate $y_n'$ from $G''/w$ using
$G'/w=y_n'+Sy_n$; then
$p(S'-S^2)=q'-p''-q(q-p')/p$. The other two follow by substitution.
The opposite signs printed in (12.14) and (12.15) would cancel each
other, but neither intermediate identity is valid with those signs.

\Needspace{5\baselineskip}
\par\bigskip
\noindent\textbf{12.1.5}\quad Given the following integral representations, in
which \erratum{$n\in\mathbb Z$} and the contours \erratum{encircle the origin once counterclockwise} but no other singular points, derive
the corresponding generating functions:
\begin{enumerate}
\item[(a)] Bessel functions:
  \[
  J_{n}(x)=\frac{1}{2\pi i}\oint e^{(x/2)(t-1/t)}t^{-n-1}\,dt.
  \]
\item[(b)] Modified Bessel functions:
  \[
  I_{n}(x)=\frac{1}{2\pi i}\oint e^{(x/2)(t+1/t)}t^{-n-1}\,dt.
  \]
\end{enumerate}
\solution
Both integrands are analytic in the punctured $t$-plane, so their
Laurent coefficients on $0<|t|<\infty$ are extracted by the given
positively oriented contour integrals. Uniqueness of the Laurent
expansion therefore gives
\[
\boxed{e^{(x/2)(t-1/t)}=\sum_{n=-\infty}^{\infty}J_n(x)t^n},
\qquad
\boxed{e^{(x/2)(t+1/t)}=\sum_{n=-\infty}^{\infty}I_n(x)t^n}.
\]
For $x\ne0$, the essential singularity at zero requires the full integer
range rather than a Taylor series. At $x=0$ both generating functions
are one, so only the $n=0$ coefficient survives.
To exhibit the coefficients directly, multiply the two exponential series:
\[
e^{xt/2}e^{\mp x/(2t)}
=\sum_{j,k\ge0}\frac{(x/2)^{j+k}(\mp1)^k}{j!k!}t^{j-k}.
\]
For $n\ge0$, the terms with $j-k=n$ have $j=k+n$ and give
\[
J_n(x)=\sum_{k\ge0}\frac{(-1)^k(x/2)^{n+2k}}{k!(n+k)!},\quad
I_n(x)=\sum_{k\ge0}\frac{(x/2)^{n+2k}}{k!(n+k)!}.
\]
The negative coefficients similarly give $J_{-n}=(-1)^nJ_n$ and
$I_{-n}=I_n$. On any compact annulus in $t$, the product of exponential
series converges absolutely and uniformly, justifying collection of
coefficients and contour integration term by term.

\Needspace{5\baselineskip}
\par\bigskip
\noindent\textbf{12.1.6}\quad Expand the generating function for the Legendre
polynomials, $(1-2tz+t^{2})^{-1/2}$, in powers of $t$. Assume that $t$ is
small. Collect the coefficients of $t^{0}$, $t^{1}$, and $t^{2}$.
\par\nobreak\hfill\emph{ANS.}\enspace
$a_{0}=P_{0}(z)=1$,\quad$a_{1}=P_{1}(z)=z$,\quad
$a_{2}=P_{2}(z)=\frac{1}{2}(3z^{2}-1)$.
\solution
Use $(1+u)^{-1/2}=1-u/2+3u^2/8+O(u^3)$ with
$u=-2tz+t^2$. To retain only powers through $t^2$,
\[
-\frac u2=zt-\frac{t^2}2,\qquad
\frac{3u^2}{8}=\frac38(4z^2t^2+O(t^3))
=\frac32z^2t^2+O(t^3).
\]
Consequently
\[
(1-2tz+t^2)^{-1/2}
=1+zt+\frac{3z^2-1}{2}t^2+O(t^3).
\]
Thus $\boxed{P_0=1,\ P_1=z,\ P_2=(3z^2-1)/2}$.
Here $z$ is fixed and the branch is chosen to equal one at $t=0$;
the expansion is valid in a sufficiently small disk about zero.

\Needspace{5\baselineskip}
\par\bigskip
\noindent\textbf{12.1.7}\quad The set of Chebyshev polynomials usually denoted
$U_{n}(x)$ has the generating-function formula
\[
\frac{1}{1-2xt+t^{2}}=\sum_{n=0}^{\infty}U_{n}(x)t^{n}.
\]
Derive a recurrence formula (for integer $n\ge0$) connecting three $U_{n}$ of
consecutive $n$, \erratum{using $U_{-1}(x)=0$ at $n=0$}.
\solution
Multiply the generating-function identity by its denominator:
\[
(1-2xt+t^2)\sum_{n=0}^{\infty}U_n(x)t^n=1.
\]
The constant term gives $U_0=1$, and the coefficient of $t$ gives
$U_1-2xU_0=0$, so $U_1=2x$. At power $t^{n+1}$ for $n\ge1$,
\[
U_{n+1}-2xU_n+U_{n-1}=0.
\]
Hence
\[
\boxed{U_{n+1}=2xU_n-U_{n-1}\quad(n\ge1)}.
\]
\erratum{To write this recurrence at $n=0$ as requested, define
$U_{-1}=0$; otherwise the three-term formula starts at $n=1$ and the
$n=0$ case is the separate initial equation $U_1=2xU_0$.}
The recurrence gives $U_2=4x^2-1$ and $U_3=8x^3-4x$, which also
follow by direct expansion of the rational generating function.

\Needspace{8\baselineskip}
\subsection*{12.2 Bernoulli Numbers}

\Needspace{5\baselineskip}
\par\bigskip
\noindent\textbf{12.2.1}\quad Verify the identities, Eqs.~(12.32) and (12.46).
\solution
For Eq. (12.32), multiply the numerator and denominator involving
$e^{-t}$ by $e^t$:
\[
\frac{-t}{e^{-t}-1}=\frac{te^t}{e^t-1}
=\frac{t(e^t-1)+t}{e^t-1}=t+\frac{t}{e^t-1}.
\]
Subtracting $t$ gives
\[
\boxed{\frac{t}{e^t-1}=\frac{-t}{e^{-t}-1}-t}.
\]
The first equality already proves Eq. (12.46):
\[
\boxed{\frac{te^t}{e^t-1}=\frac{-t}{e^{-t}-1}}.
\]
These identities hold away from the denominator's zeros and extend
through the removable point $t=0$ by continuity; both quotients there
have limiting value one.

\Needspace{5\baselineskip}
\par\bigskip
\noindent\textbf{12.2.2}\quad Show that the first Bernoulli polynomials are
\begin{align*}
B_{0}(s)&=1\\
B_{1}(s)&=s-\tfrac{1}{2}\\
B_{2}(s)&=s^{2}-s+\tfrac{1}{6}.
\end{align*}
Note that $B_{n}(0)=B_{n}$, the Bernoulli number.
\solution
The defining generating function is
$te^{st}/(e^t-1)=\sum_{n\ge0}B_n(s)t^n/n!$.
First invert the series $(e^t-1)/t=1+t/2+t^2/6+O(t^3)$.
If its inverse is $1+at+bt^2+O(t^3)$, the product requires
$a+1/2=0$ and $b+a/2+1/6=0$. Thus $a=-1/2$, $b=1/12$.
Multiply by $e^{st}=1+st+s^2t^2/2+O(t^3)$:
\[
\frac{te^{st}}{e^t-1}
=1+\left(s-\frac12\right)t
+\left(\frac{s^2}{2}-\frac s2+\frac1{12}\right)t^2+O(t^3).
\]
Comparing with $B_0+B_1t+B_2t^2/2!+\cdots$ yields
\[
\boxed{B_0(s)=1,\quad B_1(s)=s-\frac12,\quad
 B_2(s)=s^2-s+\frac16}.
\]
At $s=0$ the generating function is the Bernoulli-number generating
function itself, proving $B_n(0)=B_n$ for every $n$.

\Needspace{5\baselineskip}
\par\bigskip
\noindent\textbf{12.2.3}\quad Show that
\[
\tan x=\sum_{n=1}^{\infty}
\frac{(-1)^{n-1}2^{2n}(2^{2n}-1)B_{2n}}{(2n)!}x^{2n-1},
\quad-\frac{\pi}{2}<x<\frac{\pi}{2}.
\]
\emph{Hint.}~$\tan x=\cot x-2\cot2x$.
\solution
First verify the hinted identity:
\[
\cot x-2\cot2x
=\frac{\cos x}{\sin x}
 -\frac{\cos^2x-\sin^2x}{\sin x\cos x}
=\frac{\sin x}{\cos x}=\tan x.
\]
The identity extends through zero by continuity after the poles cancel.
From the Bernoulli generating function,
\[
\cot x=i+\frac{2i}{e^{2ix}-1}
=\frac1x+\sum_{n=1}^{\infty}
 \frac{(-1)^n2^{2n}B_{2n}}{(2n)!}x^{2n-1},\quad |x|<\pi.
\]
Here the $B_1=-1/2$ term cancels the initial $i$, and all odd Bernoulli
numbers beyond $B_1$ vanish. Substitute this series at $x$ and at $2x$:
\begin{align*}
\cot x-2\cot2x
&=\left(\frac1x-\frac2{2x}\right)
+\sum_{n\ge1}\frac{(-1)^n2^{2n}B_{2n}}{(2n)!}
 [1-2^{2n}]x^{2n-1}\\
&=\boxed{\sum_{n=1}^{\infty}
\frac{(-1)^{n-1}2^{2n}(2^{2n}-1)B_{2n}}{(2n)!}x^{2n-1}}.
\end{align*}
Both cotangent expansions converge for $|x|<\pi/2$. The nearest poles
of $\tan x$ are at $\pm\pi/2$, so this is also its exact Taylor radius.
The first terms are $x+x^3/3+2x^5/15+\cdots$.

\Needspace{8\baselineskip}
\subsection*{12.3 Euler-Maclaurin Integration Formula}

\Needspace{5\baselineskip}
\par\bigskip
\noindent\textbf{12.3.1}\quad The Euler-Maclaurin integration formula may be
used for the evaluation of finite series:
\[
\sum_{m=1}^{n}f(m)=\int_{1}^{n}f(x)\,dx+\frac{1}{2}f(1)+\frac{1}{2}f(n)
+\frac{B_{2}}{2!}\left[f'(n)-f'(1)\right]+\cdots.
\]
Show that
\begin{enumerate}
\item[(a)] $\displaystyle\sum_{m=1}^{n}m=\tfrac{1}{2}n(n+1)$.
\item[(b)] $\displaystyle\sum_{m=1}^{n}m^{2}=\tfrac{1}{6}n(n+1)(2n+1)$.
\item[(c)] $\displaystyle\sum_{m=1}^{n}m^{3}=\tfrac{1}{4}n^{2}(n+1)^{2}$.
\item[(d)] $\displaystyle\sum_{m=1}^{n}m^{4}
  =\tfrac{1}{30}n(n+1)(2n+1)(3n^{2}+3n-1)$.
\end{enumerate}
\solution
For each polynomial, the Euler--Maclaurin correction sequence terminates
because sufficiently high derivatives vanish. Use $B_2=1/6$, $B_4=-1/30$.

(a) For $f(x)=x$, the first-derivative difference is $1-1=0$, so
\[
\sum_{m=1}^nm=\frac{n^2-1}{2}+\frac{1+n}{2}
=\boxed{\frac{n(n+1)}2}.
\]
(b) For $f(x)=x^2$, $f'=2x$ and higher odd derivatives vanish. Hence
\begin{align*}
\sum_{m=1}^nm^2
&=\frac{n^3-1}{3}+\frac{1+n^2}{2}+\frac1{12}(2n-2)\\
&=\frac{2n^3+3n^2+n}{6}
=\boxed{\frac{n(n+1)(2n+1)}6}.
\end{align*}
(c) For $f(x)=x^3$, $f'=3x^2$ and $f^{(3)}=6$, whose endpoint
difference is zero. Thus
\begin{align*}
\sum_{m=1}^nm^3
&=\frac{n^4-1}{4}+\frac{1+n^3}{2}+\frac1{12}(3n^2-3)\\
&=\frac{n^4+2n^3+n^2}{4}
=\boxed{\frac{n^2(n+1)^2}{4}}.
\end{align*}
(d) For $f(x)=x^4$, $f'=4x^3$, $f^{(3)}=24x$, and $f^{(5)}=0$.
Retain both nonzero endpoint corrections:
\begin{align*}
\sum_{m=1}^nm^4
&=\frac{n^5-1}{5}+\frac{1+n^4}{2}
 +\frac1{12}(4n^3-4)-\frac1{720}(24n-24)\\
&=\frac{6n^5+15n^4+10n^3-n}{30}\\
&=\boxed{\frac{n(n+1)(2n+1)(3n^2+3n-1)}{30}}.
\end{align*}
The last factorization follows from
$(n+1)(2n+1)(3n^2+3n-1)=6n^4+15n^3+10n^2-1$.

\Needspace{5\baselineskip}
\par\bigskip
\noindent\textbf{12.3.2}\quad The Euler-Maclaurin integration formula provides
a way of calculating the Euler-Mascheroni constant $\gamma$ to high accuracy.
Using $f(x)=1/x$ in Eq.~(12.57) (with interval $[1,n]$) and the definition of
$\gamma$, Eq.~(1.13), we obtain
\[
\gamma=\sum_{s=1}^{n}s^{-1}-\ln n-\frac{1}{2n}
+\sum_{k=1}^{N}\frac{B_{2k}}{(2k)n^{2k}}\erratum{+R_{n,N}}.
\]
Using double-precision arithmetic, calculate the approximation for
$N=1,2,\ldots$ \erratum{at $n=1000$}, as in the printed answer below.

\emph{Note.}~See D.~E. Knuth, Euler's constant to 1271 places.
\emph{Math.\ Comput.}~\textbf{16}: 275 (1962).
\par\nobreak\hfill\emph{ANS.}\enspace
For $n=1000$, $N=2$\quad$\gamma\erratum{\approx}0.5772\,1566\,4901$
\erratum{(truncated, rather than rounded, to twelve decimal places)}.
\solution
\erratum{The printed finite-$N$ equality omits the remainder $R_{n,N}$; it is restored above.}
Put $H_n=\sum_{s=1}^n1/s$ and define the computable approximation
\[
A_{n,N}=H_n-\ln n-\frac1{2n}
 +\sum_{k=1}^N\frac{B_{2k}}{2k\,n^{2k}}.
\]
To see the signs, apply Euler--Maclaurin to the tail from $n$ to $M$.
Since $f^{(2k-1)}(x)=-(2k-1)!/x^{2k}$, it gives
\[
H_M-H_n=\ln(M/n)+\frac1{2M}-\frac1{2n}
 +\sum_{k=1}^N\frac{B_{2k}}{2k}
 (n^{-2k}-M^{-2k})+\text{remainder}.
\]
Let $M\to\infty$ and use $H_M-\ln M\to\gamma$.
This yields $\gamma=A_{n,N}+R_{n,N}$. A rigorous signed remainder
bound, derived in Exercise 12.6.6, is
\[
0<(-1)^N R_{n,N}
<\frac{|B_{2N+2}|}{(2N+2)n^{2N+2}}.
\]

Using IEEE double arithmetic, $n=1000$, an accurately summed harmonic
sum, and then accurately summing the corrections, gives
\begin{center}
\begin{tabular}{rrr}\toprule
$N$&$A_{1000,N}$&last correction\\\midrule
1&$0.5772156649015416$&$8.3333333333\times10^{-8}$\\
2&$0.5772156649015333$&$-8.3333333333\times10^{-15}$\\
3&$0.5772156649015333$&$3.9682539683\times10^{-21}$\\
4&$0.5772156649015333$&$-4.1666666667\times10^{-27}$\\
5&$0.5772156649015333$&$7.5757575758\times10^{-33}$\\
6&$0.5772156649015333$&$-2.1092796093\times10^{-38}$\\
\bottomrule\end{tabular}\end{center}
The exact truncation error at $N=2$ is below
$1/(252\,1000^6)=3.9682539683\times10^{-21}$.
The remaining roughly $4\times10^{-16}$ difference from
$\gamma=0.5772156649015328606\ldots$ is double-arithmetic rounding
in the harmonic sum and subtraction. More corrections cannot recover
those already lost digits. To twelve decimal places the result is
$\boxed{\gamma=0.577215664902}$, consistent with the textbook's
truncated display $0.5772\,1566\,4901$.

\Needspace{8\baselineskip}
\subsection*{12.4 Dirichlet Series}

\Needspace{5\baselineskip}
\par\bigskip
\noindent\textbf{12.4.1}\quad From
$B_{2n}=(-1)^{n-1}\dfrac{2(2n)!}{(2\pi)^{2n}}\zeta(2n)$, show that
\begin{enumerate}
\item[(a)] $\zeta(2)=\dfrac{\pi^{2}}{6}$,\qquad
(d) $\zeta(8)=\dfrac{\pi^{8}}{9450}$,
\item[(b)] $\zeta(4)=\dfrac{\pi^{4}}{90}$,\qquad
(e) $\zeta(10)=\dfrac{\pi^{10}}{93{,}555}$,
\item[(c)] $\zeta(6)=\dfrac{\pi^{6}}{945}$,
\end{enumerate}
\solution
Rearrange the stated identity before inserting numbers:
\[
\zeta(2n)=\frac{(-1)^{n-1}B_{2n}(2\pi)^{2n}}{2(2n)!}.
\]
With $B_2=1/6$, $B_4=-1/30$, $B_6=1/42$, $B_8=-1/30$,
$B_{10}=5/66$, the five evaluations are
\begin{align*}
\zeta(2)&=\frac{(1/6)4\pi^2}{2\cdot2}=\boxed{\pi^2/6},\\
\zeta(4)&=\frac{(1/30)16\pi^4}{2\cdot24}=\boxed{\pi^4/90},\\
\zeta(6)&=\frac{(1/42)64\pi^6}{2\cdot720}=\boxed{\pi^6/945},\\
\zeta(8)&=\frac{(1/30)256\pi^8}{2\cdot40320}=\boxed{\pi^8/9450},\\
\zeta(10)&=\frac{(5/66)1024\pi^{10}}{2\cdot3628800}
=\boxed{\pi^{10}/93555}.
\end{align*}
Each alternating sign cancels the sign of the corresponding Bernoulli
number, leaving the positive values required by the defining zeta series.

\Needspace{5\baselineskip}
\par\bigskip
\noindent\textbf{12.4.2}\quad The integral
\[
\int_{0}^{1}\frac{dx}{x}\,[\ln(1-x)]^{2}
\]
appears in the fourth-order correction to the magnetic moment of the
electron. Show that it equals $2\zeta(3)$.

\emph{Hint.}~Let $1-x=e^{-t}$.
\solution
Let $1-x=e^{-t}$, so $x=1-e^{-t}$, $dx=e^{-t}dt$, and
$\ln(1-x)=-t$. The endpoints $x=0,1$ become $t=0,\infty$, giving
\[
I=\int_0^\infty\frac{t^2e^{-t}}{1-e^{-t}}dt.
\]
For $t>0$, $e^{-t}/(1-e^{-t})=\sum_{m=1}^\infty e^{-mt}$.
All summands are nonnegative, so Tonelli's theorem permits interchange:
\[
I=\sum_{m=1}^\infty\int_0^\infty t^2e^{-mt}dt.
\]
Put $u=mt$ in each integral. Two integrations by parts give
$\int_0^\infty u^2e^{-u}du=2\int_0^\infty ue^{-u}du=2$.
Consequently
\[
\boxed{I=2\sum_{m=1}^\infty\frac1{m^3}=2\zeta(3)}.
\]
The finite value also proves convergence of the original improper
integral, including its logarithmic endpoint singularity.

\Needspace{5\baselineskip}
\par\bigskip
\noindent\textbf{12.4.3}\quad
\begin{enumerate}
\item[(a)] Show that
  \[
  \int_{0}^{\infty}\frac{(\ln z)^{2}}{1+z^{2}}\,dz
  =4\left(1-\frac{1}{3^{3}}+\frac{1}{5^{3}}-\frac{1}{7^{3}}+\cdots\right).
  \]
\item[(b)] By contour integration show that \erratum{the integral, equivalently
  four times the parenthesized series,} evaluates to $\pi^{3}/8$.
\end{enumerate}
\solution
(a) In the part from $1$ to $\infty$, set $z=1/t$. Since
$dz/(1+z^2)=-dt/(1+t^2)$ and $(\ln z)^2=(\ln t)^2$,
\[
I=2\int_0^1\frac{(\ln t)^2}{1+t^2}dt.
\]
Expand the denominator geometrically. The absolute sum of the integrals
is finite, because
\[
\int_0^1t^{2k}(\ln t)^2dt
=\int_0^\infty u^2e^{-(2k+1)u}du=\frac2{(2k+1)^3}.
\]
Here $t=e^{-u}$ was used. Thus termwise integration gives
\[
\boxed{I=4\sum_{k=0}^\infty\frac{(-1)^k}{(2k+1)^3}=4\beta(3)}.
\]

(b) Integrate $F(z)=(\operatorname{Log}z)^2/(1+z^2)$ round an upper
semicircle, indenting the origin, using $0<\arg z<\pi$ inside the
contour. The large arc is $O((\ln R)^2/R)$ and the small arc is
$O(\epsilon|\ln\epsilon|^2)$; both vanish. The sole enclosed pole is
$z=i$, where
\[
\operatorname{Res}_{z=i}F
=\frac{(i\pi/2)^2}{2i}=\frac{i\pi^2}{8},\qquad
\oint Fdz=-\frac{\pi^3}{4}.
\]
On the negative real edge, $z=-t$ and $\operatorname{Log}z=\ln t+i\pi$.
The two real-axis contributions therefore add to
\[
2I+2i\pi J-\pi^2K,\qquad
J=\int_0^\infty\frac{\ln t}{1+t^2}dt,\quad
K=\int_0^\infty\frac{dt}{1+t^2}=\frac\pi2.
\]
The substitution $t=1/u$ gives $J=-J$, so $J=0$. Equating to the
residue result yields $2I-\pi^3/2=-\pi^3/4$, hence
\[
\boxed{I=\pi^3/8},\qquad \boxed{\beta(3)=\pi^3/32}.
\]
\erratum{The value $\pi^3/8$ refers to the integral, or to four times
the series in parentheses in (a). The unmultiplied alternating odd-cube
series has value $\pi^3/32$.}

\Needspace{5\baselineskip}
\par\bigskip
\noindent\textbf{12.4.4}\quad Show that Catalan's constant, $\beta(2)$, may be
written as
\[
\beta(2)=2\sum_{k=1}^{\infty}(4k-3)^{-2}-\frac{\pi^{2}}{8}.
\]
\emph{Hint.}~$\pi^{2}=6\zeta(2)$.
\solution
Separate the positive odd integers into their two residue classes:
\[
A=\sum_{k\ge1}(4k-3)^{-2},\qquad
B=\sum_{k\ge1}(4k-1)^{-2}.
\]
Absolute convergence permits grouping. By definition $\beta(2)=A-B$,
while
\[
A+B=\sum_{j\ge0}(2j+1)^{-2}
=\zeta(2)-\sum_{j\ge1}(2j)^{-2}
=\frac34\zeta(2)=\frac{\pi^2}{8}.
\]
Eliminate $B=\pi^2/8-A$ to obtain
\[
\boxed{\beta(2)=2\sum_{k=1}^\infty(4k-3)^{-2}-\frac{\pi^2}{8}}.
\]

\Needspace{5\baselineskip}
\par\bigskip
\noindent\textbf{12.4.5}\quad Show that
\begin{enumerate}
\item[(a)] $\displaystyle\int_{0}^{1}\frac{\ln(1+x)}{x}\,dx=\frac{1}{2}\zeta(2)$,\quad
(b) $\displaystyle\lim_{a\to1^{-}}\int_{0}^{a}\frac{\ln(1-x)}{x}\,dx=\erratum{-}\zeta(2)$.
\end{enumerate}
Note that the integrand in part (b) diverges for $a=1$ but that the integral
is convergent.
\solution
(a) For $0\le x<1$, expand
$\ln(1+x)/x=\sum_{k\ge1}(-1)^{k+1}x^{k-1}/k$.
The sum of the integrals of the absolute values is
$\sum_{k\ge1}1/k^2<\infty$, which justifies termwise integration up
to the endpoint. Therefore
\[
\int_0^1\frac{\ln(1+x)}x dx
=\sum_{k\ge1}\frac{(-1)^{k+1}}{k^2}
=\zeta(2)-2\sum_{j\ge1}\frac1{(2j)^2}
=\boxed{\frac12\zeta(2)}.
\]
(b) For $a<1$, the logarithm series is uniformly convergent on
$[0,a]$ after removing the harmless limit at zero, giving
\[
\int_0^a\frac{\ln(1-x)}x dx
=-\sum_{k=1}^\infty\frac{a^k}{k^2}.
\]
Dominated convergence, with dominating summable sequence $1/k^2$,
permits $a\uparrow1$, so
\[
\boxed{\lim_{a\to1^-}\int_0^a\frac{\ln(1-x)}x dx=-\zeta(2)}.
\]
\erratum{The printed right-hand side in (b) has the wrong sign.
The integral is negative since $\ln(1-x)/x<0$ for $0<x<1$.}
The integrand diverges logarithmically at one, but its absolute integral
is $\zeta(2)$ and is finite.

\Needspace{5\baselineskip}
\par\bigskip
\noindent\textbf{12.4.6}\quad
\begin{enumerate}
\item[(a)] Show that the equation
  $\ln2=\sum_{s=1}^{\infty}(-1)^{s+1}s^{-1}$, Eq.~(1.53), may be rewritten as
  \[
  \ln2=\sum_{s=2}^{n}2^{-s}\zeta(s)
  +\sum_{p=1}^{\infty}(2p)^{-n-1}\left[1-\frac{1}{2p}\right]^{-1}.
  \]
  \emph{Hint.}~Take the terms in pairs.
\item[(b)] Calculate $\ln2$ to six significant figures.
\end{enumerate}
\solution
(a) Pair neighboring terms of the convergent alternating harmonic series:
\[
\ln2=\sum_{p=1}^\infty\left(\frac1{2p-1}-\frac1{2p}\right)
=\sum_{p=1}^\infty\frac{(2p)^{-2}}{1-(2p)^{-1}}.
\]
For $r=(2p)^{-1}$ the finite geometric identity is
\[
\frac{r^2}{1-r}=\sum_{s=2}^{n}r^s+\frac{r^{n+1}}{1-r}.
\]
Sum it over $p$. The first sum has finitely many $s$, and all terms
are nonnegative, giving exactly
\[
\boxed{\ln2=\sum_{s=2}^{n}2^{-s}\zeta(s)
+\sum_{p=1}^\infty\frac{(2p)^{-n-1}}{1-(2p)^{-1}}}.
\]
The usual nonempty first sum takes $n\ge2$; the $n=1$ case also holds
if that sum is interpreted as empty.

(b) The remainder $R_n$ is positive and, since
$1-(2p)^{-1}\ge1/2$, obeys
\[
0<R_n\le2\sum_{p\ge1}(2p)^{-n-1}=2^{-n}\zeta(n+1).
\]
Take $n=24$. Direct summation of the 23 zeta terms gives
\[
S_{24}=0.6931471209552993497\ldots,\qquad
R_{24}<5.9604646552\times10^{-8}.
\]
Both ends of the interval $S_{24}<\ln2<S_{24}+R_{24}^{\rm bound}$
round to the same six-significant-figure value:
\[
\boxed{\ln2=0.693147\quad\text{(six significant figures)}}.
\]
The error bound comes from the derived remainder, independently of any
precomputed logarithm value.

\Needspace{5\baselineskip}
\par\bigskip
\noindent\textbf{12.4.7}\quad
\begin{enumerate}
\item[(a)] Show that the equation
  $\pi/4=\sum_{s=1}^{\infty}(-1)^{s+1}(2s-1)^{-1}$, Eq.~(12.72), may be
  rewritten as
  \[
  \frac{\pi}{4}=1-2\sum_{s=1}^{n}4^{-2s}\zeta(2s)
  -2\sum_{p=1}^{\infty}(4p)^{-2n-2}\left[1-\frac{1}{(4p)^{2}}\right]^{-1}.
  \]
\item[(b)] Calculate $\pi/4$ to six significant figures.
\end{enumerate}
\solution
(a) Separate the first term of the arctangent series, then pair the rest:
\[
\frac\pi4=1+\sum_{p=1}^\infty
\left(-\frac1{4p-1}+\frac1{4p+1}\right)
=1-2\sum_{p=1}^\infty\frac1{(4p)^2-1}.
\]
With $r=(4p)^{-2}$,
\[
\frac1{(4p)^2-1}=\frac r{1-r}
=\sum_{s=1}^nr^s+\frac{r^{n+1}}{1-r}.
\]
Summing over $p$ yields
\[
\boxed{\frac\pi4=1-2\sum_{s=1}^n4^{-2s}\zeta(2s)
-2\sum_{p=1}^\infty
\frac{(4p)^{-2n-2}}{1-(4p)^{-2}}}.
\]
(b) The omitted correction has negative sign. Its magnitude is bounded
using $1-(4p)^{-2}\ge15/16$:
\[
0<R_n:=2\sum_{p\ge1}\frac{(4p)^{-2n-2}}{1-(4p)^{-2}}
\le\frac{32}{15}\,4^{-2n-2}\zeta(2n+2).
\]
For $n=6$, the finite part and bound are
\[
S_6=0.7853981713451978431\ldots,\qquad
R_6<7.9477727265\times10^{-9}.
\]
Thus $S_6-R_6^{\rm bound}<\pi/4<S_6$, whose endpoints both round to
\[
\boxed{\pi/4=0.785398\quad\text{(six significant figures)}}.
\]

\Needspace{8\baselineskip}
\subsection*{12.5 Infinite Products}

\Needspace{5\baselineskip}
\par\bigskip
\noindent\textbf{12.5.1}\quad \erratum{Assume $a_n\ge0$, no factor vanishes, and the factors
are eventually positive. Product convergence means a finite nonzero
limit.} Using
\[
\ln\prod_{n=1}^{\infty}(1\pm a_{n})=\sum_{n=1}^{\infty}\ln(1\pm a_{n})
\]
and the Maclaurin expansion of $\ln(1\pm a_{n})$, show that the infinite
product $\prod_{n=1}^{\infty}(1\pm a_{n})$ converges or diverges with the
infinite series $\sum_{n=1}^{\infty}a_{n}$.
\solution
\erratum{The unqualified statement is false for signed or complex $a_n$.
It is valid in its usual comparison-test form when $a_n\ge0$ and the
factors are eventually positive. A convergent infinite product here means
a finite, nonzero limit, with no zero factor.}
Discarding finitely many initial nonzero factors only multiplies all later
partial products by a fixed nonzero constant. We therefore apply real
logarithms to the eventually positive tail. For that tail, convergence
of the product is equivalent to convergence of its logarithmic series. If $a_n\not\to0$, a nonzero product limit
is already impossible, since the ratio of two successive partial products
would not tend to one. Otherwise, eventually $0\le a_n\le1/2$.
Integration of $1/(1+t)$ and $1/(1-t)$ on $[0,a_n]$ gives
\[
\frac{a_n}{2}\le\ln(1+a_n)\le a_n,\qquad
 a_n\le-\ln(1-a_n)\le2a_n.
\]
The comparison test now proves the stated equivalence for both choices
of sign. If $\sum a_n$ diverges, the plus-product grows without bound
and the minus-product tends to zero.

To see why the hypotheses matter, take
$a_n=(-1)^n/\sqrt{n+1}$. Its series converges by the alternating test,
and every $1+a_n$ is positive. But the logarithm expansion gives
\[
\ln(1+a_n)=a_n-\frac12a_n^2+O(|a_n|^3).
\]
Here $\sum a_n^2=\sum1/(n+1)=\infty$, while
$\sum|a_n|^3<\infty$. Consequently the logarithmic partial sums tend
to $-\infty$, and $\prod(1+a_n)$ tends to zero, not a nonzero limit.

A useful corrected version for signed or complex terms is: if
$a_n\to0$, $\sum|a_n|^2<\infty$, and no factor vanishes, then
$\prod(1+a_n)$ converges to a nonzero limit exactly when $\sum a_n$
converges. Indeed $\log(1+a_n)-a_n=O(|a_n|^2)$, using the logarithm
branch near one; its error series is absolutely convergent. To justify
necessity as well as sufficiency, suppose the partial products $P_n\to L\ne0$.
Choose a disk about $L$ avoiding zero and one analytic logarithm $\ell$
there. For large $n$, $\ell(P_n)-\ell(P_{n-1})=\log(1+a_n)$:
both sides tend to zero, and their difference is an integer multiple
of $2\pi i$. The logarithmic tail telescopes to a convergent limit.
Conversely, a convergent logarithmic series exponentiates to a nonzero
product. Subtracting the absolutely convergent error proves the claim
about $\sum a_n$. The minus case follows by replacing $a_n$ by $-a_n$.

\Needspace{5\baselineskip}
\par\bigskip
\noindent\textbf{12.5.2}\quad An infinite product appears in the form
\[
\prod_{n=1}^{\infty}\left(\frac{1+a/n}{1+b/n}\right),
\]
where $a$ and $b$ are constants \erratum{not equal to negative integers,
so all factors are defined and nonzero}. Show that this infinite product
converges \erratum{to a finite nonzero limit} only if $a=b$.
\solution
\erratum{Use the standard finite, nonzero meaning of product convergence,
and assume every denominator and numerator factor is nonzero. If a zero
limit were counted as convergence, the asserted ``only if'' would fail.}
For fixed real or complex $a,b$, the tail has a logarithm branch near one.
Its expansion is
\[
\log\frac{1+a/n}{1+b/n}
=\frac{a-b}{n}+O(n^{-2}).
\]
The error series converges absolutely, so for the partial product $P_N$,
\[
\log P_N=(a-b)H_N+C+o(1)
=(a-b)\log N+C_1+o(1).
\]
The finitely many initial factors only alter the nonzero multiplicative
constant. Thus $P_N=K N^{a-b}(1+o(1))$, $K\ne0$.
If $\Re(a-b)>0$, its modulus tends to infinity; if $\Re(a-b)<0$,
it tends to zero. If $a-b=i\beta$ with real $\beta\ne0$, its phase
$\beta\ln N$ has no limit modulo $2\pi$: integers nearest
$e^{2\pi k/|\beta|}$ and $e^{(2\pi k+\pi)/|\beta|}$ give opposite
limiting phases. Hence no finite nonzero limit exists unless $a=b$.
When $a=b$, every defined factor equals one, so
\[
\boxed{\text{the nonzero product converges exactly when }a=b.}
\]
For a concrete zero-limit exception, $a=0,b=1$ gives
$P_N=\prod_{n=1}^N n/(n+1)=1/(N+1)\to0$.
If $a$ is a negative integer a numerator factor is zero; if $b$ is a
negative integer the written product is undefined at that factor.

\Needspace{5\baselineskip}
\par\bigskip
\noindent\textbf{12.5.3}\quad Show that the infinite product representations
of $\sin x$ and $\cos x$ are consistent with the identity
$2\sin x\cos x=\sin2x$.
\solution
The product formulas are
\[
\sin x=x\prod_{n=1}^{\infty}\left(1-\frac{x^2}{n^2\pi^2}\right),\quad
\cos x=\prod_{n=1}^{\infty}
 \left(1-\frac{4x^2}{(2n-1)^2\pi^2}\right).
\]
In the sine product rewrite $x^2/(n^2\pi^2)=4x^2/[(2n)^2\pi^2]$.
The two products then supply the even and odd positive integers:
\begin{align*}
2\sin x\cos x
&=2x\prod_{n\ge1}\left(1-\frac{4x^2}{(2n)^2\pi^2}\right)
       \prod_{n\ge1}\left(1-\frac{4x^2}{(2n-1)^2\pi^2}\right)\\
&=2x\prod_{m=1}^{\infty}\left(1-\frac{(2x)^2}{m^2\pi^2}\right)
=\boxed{\sin2x}.
\end{align*}
The deviations of the factors from one are bounded on compact $x$ sets
by a constant times $1/n^2$. Their absolute summability justifies
interleaving the factors. At zeros of a factor the equality extends
by continuity, so it holds for all complex $x$.

\Needspace{5\baselineskip}
\par\bigskip
\noindent\textbf{12.5.4}\quad Determine the limit to which
$\displaystyle\prod_{n=2}^{\infty}\left(1+\frac{(-1)^{n}}{n}\right)$
converges.
\solution
Pair the factor with even index $2k$ with the following odd index:
\[
\left(1+\frac1{2k}\right)
\left(1-\frac1{2k+1}\right)
=\frac{2k+1}{2k}\frac{2k}{2k+1}=1.
\]
To ensure that pairing has not hidden a subsequence issue, let
$P_N=\prod_{n=2}^N[1+(-1)^n/n]$. Then exactly
\[
P_{2M+1}=1,\qquad P_{2M}=\frac{2M+1}{2M}.
\]
Both subsequences tend to one. Therefore the complete sequence of
partial products converges and
\[
\boxed{\prod_{n=2}^{\infty}\left(1+\frac{(-1)^n}{n}\right)=1}.
\]

\Needspace{5\baselineskip}
\par\bigskip
\noindent\textbf{12.5.5}\quad Show that
$\displaystyle\prod_{n=2}^{\infty}\left[1-\frac{2}{n(n+1)}\right]
=\frac{1}{3}$.
\solution
Factor each term before taking any limit:
\[
1-\frac2{n(n+1)}
=\frac{n^2+n-2}{n(n+1)}
=\frac{(n-1)(n+2)}{n(n+1)}.
\]
The finite product telescopes in two separate pieces:
\begin{align*}
P_N&=\left(\prod_{n=2}^N\frac{n-1}{n}\right)
     \left(\prod_{n=2}^N\frac{n+2}{n+1}\right)\\
&=\left(\frac12\frac23\cdots\frac{N-1}{N}\right)
  \left(\frac43\frac54\cdots\frac{N+2}{N+1}\right)\\
&=\frac1N\frac{N+2}{3}.
\end{align*}
Taking $N\to\infty$ proves $\boxed{P_\infty=1/3}$.

\Needspace{5\baselineskip}
\par\bigskip
\noindent\textbf{12.5.6}\quad Prove that
$\displaystyle\prod_{n=2}^{\infty}\left(1-\frac{1}{n^{2}}\right)=\frac{1}{2}$.
\solution
Use $1-1/n^2=(n-1)(n+1)/n^2$. For a finite upper limit,
\[
P_N=\left(\prod_{n=2}^N\frac{n-1}{n}\right)
    \left(\prod_{n=2}^N\frac{n+1}{n}\right)
=\frac1N\frac{N+1}{2}=\frac{N+1}{2N}.
\]
Both telescoping factors are exact for every $N\ge2$.
Thus $\boxed{\lim_{N\to\infty}P_N=1/2}$.

\Needspace{5\baselineskip}
\par\bigskip
\noindent\textbf{12.5.7}\quad Verify the Euler identity
$\displaystyle\prod_{p=1}^{\infty}(1+z^{p})
=\prod_{q=1}^{\infty}(1-z^{2q-1})^{-1}$, $|z| < 1$.
\solution
For $|z|<1$, all factors $1\pm z^n$ are nonzero and
$\sum_{n\ge1}|z|^n<\infty$. The corresponding products and their
logarithmic tails converge absolutely, permitting the following
factorization and regrouping:
\[
\prod_{p=1}^{\infty}(1+z^p)
=\frac{\prod_{p=1}^{\infty}(1-z^{2p})}
       {\prod_{p=1}^{\infty}(1-z^p)}.
\]
Split the denominator into even and odd indices:
\[
\prod_{p\ge1}(1-z^p)
=\prod_{q\ge1}(1-z^{2q})\prod_{q\ge1}(1-z^{2q-1}).
\]
The even-index product is nonzero, so it cancels legitimately, leaving
\[
\boxed{\prod_{p=1}^{\infty}(1+z^p)
=\prod_{q=1}^{\infty}(1-z^{2q-1})^{-1}}.
\]
The same argument is uniform on $|z|\le r<1$. No boundary claim at
$|z|=1$ follows from it.

\Needspace{5\baselineskip}
\par\bigskip
\noindent\textbf{12.5.8}\quad Show that
$\prod_{r=1}^{\infty}(1+x/r)e^{-x/r}$ converges for all finite $x$ (except for
the zeros of $1+x/r$).

\emph{Hint.}~Write the $n$th factor as $1+a_{n}$.
\erratum{Here convergence means a nonzero product; at a negative integer
$x$ the partial products instead have the zero limit.}
\solution
Write the $r$th factor as $1+a_r$ with $u=x/r$:
\begin{align*}
(1+u)e^{-u}
&=(1+u)\left(1-u+\frac{u^2}{2}-\frac{u^3}{6}+\cdots\right)\\
&=1-\frac{u^2}{2}+\frac{u^3}{3}+O(u^4).
\end{align*}
Hence $a_r=-x^2/(2r^2)+O(r^{-3})$ for fixed finite complex $x$,
and $\sum|a_r|<\infty$. Once $|a_r|<1/2$, the convergent logarithm
series has the bound
\[
|\log(1+a_r)|
\le\sum_{k\ge1}\frac{|a_r|^k}{k}\le2|a_r|.
\]
The logarithmic tail therefore converges absolutely, proving a finite
nonzero product whenever no factor vanishes. The exceptional points
are exactly $x=-1,-2,\ldots$, since exponentials never vanish.

The estimate is uniform for $|x|\le R$ after finitely many factors:
$|a_r|\le C_R/r^2$. Thus the product even defines an entire function
by locally uniform convergence.
\erratum{At the excluded points it still converges as a sequence of
partial products, with value zero; they are excluded only from the
convention of a nonzero convergent infinite product.}

\Needspace{5\baselineskip}
\par\bigskip
\noindent\textbf{12.5.9}\quad Derive the formula, \erratum{for $0<x<\pi$ with real logarithms},
\[
\ln\sin x=\ln x+\sum a_{n}x^{n},
\]
giving the explicit form for the coefficients $a_{n}$.

\emph{Hint.}~$d(\ln\sin x)/dx=\cot x$.
\solution
It is convenient to remove the known logarithmic singularity first:
\[
\frac d{dx}\ln\frac{\sin x}{x}=\cot x-\frac1x.
\]
The Bernoulli expansion of the right side is
\[
\cot x-\frac1x
=\sum_{k=1}^{\infty}
 \frac{(-1)^k2^{2k}B_{2k}}{(2k)!}x^{2k-1},\qquad |x|<\pi.
\]
Integrate term by term from zero, using
$\sin x/x\to1$ to fix the integration constant to zero:
\[
\boxed{\ln\frac{\sin x}{x}
=\sum_{k=1}^{\infty}
\frac{(-1)^k2^{2k}B_{2k}}{2k(2k)!}x^{2k}}.
\]
Therefore
\[
\boxed{a_{2k}=\frac{(-1)^k2^{2k}B_{2k}}{2k(2k)!}
=-\frac{\zeta(2k)}{k\pi^{2k}},\qquad a_{2k-1}=0\quad(k\ge1)}.
\]
The constant coefficient is $a_0=0$ by the normalization at zero.
In particular,
$\ln(\sin x/x)=-x^2/6-x^4/180-x^6/2835-\cdots$.
\erratum{For real logarithms the question's form $\ln\sin x=\ln x+\cdots$
uses $0<x<\pi$. For $-\pi<x<0$, write $\ln|\sin x|=\ln|x|+\cdots$.
In complex form, the quotient has an analytic logarithm on $|x|<\pi$
normalized to zero at the origin; any separate logarithms must use
consistent branches.}

\Needspace{5\baselineskip}
\par\bigskip
\noindent\textbf{12.5.10}\quad Using the infinite product representations of
$\sin z$, show that \erratum{for $|z|<\pi$}
\[
z\cot z=1-2\sum_{m,n=1}^{\infty}\left(\frac{z}{n\pi}\right)^{2m},
\]
and hence that the Bernoulli numbers are given by the formula
\[
B_{2n}=(-1)^{n-1}\frac{2(2n)!}{(2\pi)^{2n}}\zeta(2n).
\]
This is an alternate route to Eq.~(12.38).

\emph{Hint.}~The result of Exercise 12.5.9 will be helpful.
\solution
Take the logarithmic derivative of the sine product on a compact subset
avoiding its zeros. Local uniform convergence of the derivative series
follows from its $O(n^{-2})$ terms. Thus
\[
\cot z=\frac1z+\sum_{n=1}^{\infty}
 \frac{-2z}{n^2\pi^2-z^2},\qquad
z\cot z=1-2\sum_{n=1}^{\infty}\frac{z^2}{n^2\pi^2-z^2}.
\]
\erratum{The requested double power series is restricted to $|z|<\pi$;
it is not a representation for arbitrary $z$ away from the poles.}
In that disk each denominator can be expanded geometrically:
\[
\frac{z^2}{n^2\pi^2-z^2}
=\sum_{m=1}^{\infty}\left(\frac{z}{n\pi}\right)^{2m}.
\]
For $|z|\le r<\pi$ the double absolute sum is bounded by
\[
\sum_{n\ge1}\frac{r^2/(n\pi)^2}{1-r^2/(n\pi)^2}
\le\frac{r^2\zeta(2)}{\pi^2[1-r^2/\pi^2]}<\infty,
\]
so the two summations may be interchanged. Hence
\[
\boxed{z\cot z=1-2\sum_{m,n\ge1}(z/(n\pi))^{2m}}
=1-2\sum_{m\ge1}\frac{\zeta(2m)}{\pi^{2m}}z^{2m}.
\]
Independently, the Bernoulli generating function gives
\[
z\cot z=1+\sum_{m\ge1}
 \frac{(-1)^m2^{2m}B_{2m}}{(2m)!}z^{2m}.
\]
Equate the coefficients of $z^{2m}$ and solve for $B_{2m}$:
\[
\boxed{B_{2m}=(-1)^{m-1}\frac{2(2m)!}{(2\pi)^{2m}}\zeta(2m)}.
\]

\Needspace{8\baselineskip}
\subsection*{12.6 Asymptotic Series}

\Needspace{5\baselineskip}
\par\bigskip
\noindent\textbf{12.6.1}\quad Integrating by parts, develop asymptotic
expansions \erratum{as $x\to+\infty$} of the Fresnel integrals
\begin{enumerate}
\item[(a)] $\displaystyle C(x)=\int_{0}^{x}\cos\frac{\pi u^{2}}{2}\,du$,\quad
(b) $\displaystyle s(x)=\int_{0}^{x}\sin\frac{\pi u^{2}}{2}\,du$.
\end{enumerate}
These integrals appear in the analysis of a knife-edge diffraction pattern.
\solution
Take $x\to+\infty$ and write $S(x)=s(x)$ for the sine integral in the
question. Put $\theta=\pi x^2/2$. Rotating the Gaussian integration
ray through $\pi/4$ gives
\[
\int_0^\infty e^{i\pi u^2/2}du
=e^{i\pi/4}\int_0^\infty e^{-\pi t^2/2}dt
=\frac{1+i}{2}.
\]
The rotation is justified by an arc estimate: on $u=Re^{i\phi}$,
$|e^{i\pi u^2/2}|=e^{-\pi R^2\sin(2\phi)/2}$; its integral over
$0\le\phi\le\pi/4$, including the small region next to $\phi=0$,
is bounded using $\sin2\phi\ge4\phi/\pi$:
\[
R\int_0^{\pi/4}e^{-\pi R^2\sin(2\phi)/2}d\phi
\le R\int_0^\infty e^{-2R^2\phi}d\phi=\frac1{2R}.
\]
Therefore
\[
C(x)+iS(x)=\frac{1+i}{2}-T_0(x),\qquad
T_m(x)=\int_x^\infty u^{-2m}e^{i\pi u^2/2}du.
\]
Since $d(e^{i\pi u^2/2})=i\pi u e^{i\pi u^2/2}du$, integration by
parts gives
\begin{align*}
T_m&=\left[\frac{e^{i\pi u^2/2}}{i\pi u^{2m+1}}\right]_x^\infty
 +\frac{2m+1}{i\pi}\int_x^\infty u^{-2m-2}e^{i\pi u^2/2}du\\
&=\frac{i e^{i\theta}}{\pi x^{2m+1}}
 -\frac{i(2m+1)}\pi T_{m+1}.
\end{align*}
Repeated substitution therefore gives the finite identity
\[
T_0=\frac{i e^{i\theta}}{\pi x}
\sum_{k=0}^{N-1}\frac{(-i)^k(2k-1)!!}{(\pi x^2)^k}
+\frac{(-i)^N(2N-1)!!}{\pi^N}T_N,
\]
where $(-1)!!=1$. The same integration-by-parts identity, followed by
an absolute bound on $T_{N+1}$, gives
\[
|T_N|\le\frac1{\pi x^{2N+1}}
+\frac{2N+1}{\pi}\int_x^\infty u^{-2N-2}du
=\frac2{\pi x^{2N+1}}.
\]
Thus the remainder is of the order of the next inverse power for every
fixed $N$, proving asymptotic validity.

Separate the even and odd $k$ contributions by defining
\begin{align*}
f(x)&\sim\frac1{\pi x}\left[1-\frac3{\pi^2x^4}
 +\frac{105}{\pi^4x^8}-\cdots\right],\\
g(x)&\sim\frac1{\pi^2x^3}\left[1-\frac{15}{\pi^2x^4}
 +\frac{945}{\pi^4x^8}-\cdots\right].
\end{align*}
Then $T_0\sim e^{i\theta}(g+if)$, so the requested real expansions are
\[
\boxed{C(x)\sim\frac12+f(x)\sin\theta-g(x)\cos\theta},
\qquad
\boxed{S(x)\sim\frac12-f(x)\cos\theta-g(x)\sin\theta}.
\]
For example $C=1/2+\sin\theta/(\pi x)-\cos\theta/(\pi^2x^3)
+O(x^{-5})$, with the corresponding sine-integral formula obtained
from the second box. For large negative $x$, use the oddness of both
integrals rather than the positive-infinity constants $1/2$.

\Needspace{5\baselineskip}
\par\bigskip
\noindent\textbf{12.6.2}\quad Rederive the asymptotic expansions of
$\mathrm{Ci}(x)$ and $\mathrm{si}(x)$ \erratum{as $x\to+\infty$} by repeated integration by parts.

\emph{Hint.}~$\mathrm{Ci}(x)+i\,\mathrm{si}(x)
=-\displaystyle\int_{x}^{\infty}\frac{e^{it}}{t}\,dt$.
\solution
Use the definitions in the hint, so $\mathrm{si}(x)=\mathrm{Si}(x)-\pi/2$,
and take $x\to+\infty$. Set
\[
J_m(x)=\int_x^\infty\frac{e^{it}}{t^{m+1}}dt.
\]
Integration by parts with $dv=e^{it}dt$ gives
\[
J_m=\left[\frac{e^{it}}{it^{m+1}}\right]_x^\infty
 +\frac{m+1}{i}J_{m+1}
=\frac{i e^{ix}}{x^{m+1}}-i(m+1)J_{m+1}.
\]
Consequently
\[
\mathrm{Ci}(x)+i\,\mathrm{si}(x)=-J_0
=-\frac{i e^{ix}}x\sum_{k=0}^{N-1}\frac{(-i)^kk!}{x^k}
-(-i)^NN!J_N.
\]
A further integration by parts bounds the remainder integral by
\[
|J_N|\le x^{-N-1}+(N+1)\int_x^\infty t^{-N-2}dt
=2x^{-N-1}.
\]
The residual after $N$ displayed terms is therefore $O(x^{-N-1})$.
Collect the real even/odd subsequences:
\[
A(x)\sim1-\frac{2!}{x^2}+\frac{4!}{x^4}-\cdots,\qquad
B(x)\sim1-\frac{3!}{x^2}+\frac{5!}{x^4}-\cdots.
\]
Since $-J_0\sim e^{ix}[-iA/x-B/x^2]$, multiplication by
$\cos x+i\sin x$ yields
\[
\boxed{\mathrm{Ci}(x)\sim\frac{\sin x}{x}A(x)
-\frac{\cos x}{x^2}B(x)},
\]
\[
\boxed{\mathrm{si}(x)\sim-\frac{\cos x}{x}A(x)
-\frac{\sin x}{x^2}B(x)}.
\]
These are fixed-order asymptotic truncations; the factorial coefficients
do not define convergent infinite series at a fixed finite $x$.

\Needspace{5\baselineskip}
\par\bigskip
\noindent\textbf{12.6.3}\quad Derive the asymptotic expansion of the Gauss
error function \erratum{as $x\to+\infty$}
\begin{align*}
\mathrm{erf}(x)&=\frac{2}{\sqrt{\pi}}\int_{0}^{x}e^{-t^{2}}\,dt\\
&\approx1-\frac{e^{-x^{2}}}{\sqrt{\pi}\,x}\left(1-\frac{1}{2x^{2}}
+\frac{1\cdot3}{2^{2}x^{4}}-\frac{1\cdot3\cdot5}{2^{3}x^{6}}+\cdots
+(-1)^{n}\frac{(2n-1)!!}{2^{n}x^{2n}}\right).
\end{align*}
\emph{Hint.}~$\mathrm{erf}(x)=1-\mathrm{erfc}(x)
=1-\dfrac{2}{\sqrt{\pi}}\displaystyle\int_{x}^{\infty}e^{-t^{2}}\,dt$.

Normalized so that $\mathrm{erf}(\infty)=1$, this function plays an important
role in probability theory. It may be expressed in terms of the Fresnel
integrals (Exercise 12.6.1), the incomplete gamma functions (Section 13.6),
or the confluent hypergeometric functions (Section 18.5).
\solution
For $x>0$, let
\[
I_m(x)=\int_x^\infty t^{-2m}e^{-t^2}dt,\qquad
\mathrm{erfc}(x)=\frac2{\sqrt\pi}I_0(x).
\]
Use $d(e^{-t^2})=-2t e^{-t^2}dt$ in integration by parts:
\[
I_m=\left[-\frac{e^{-t^2}}{2t^{2m+1}}\right]_x^\infty
-\frac{2m+1}{2}\int_x^\infty t^{-2m-2}e^{-t^2}dt
=\frac{e^{-x^2}}{2x^{2m+1}}-\frac{2m+1}{2}I_{m+1}.
\]
After $N$ repetitions,
\[
I_0=\frac{e^{-x^2}}{2x}
\sum_{k=0}^{N-1}\frac{(-1)^k(2k-1)!!}{2^kx^{2k}}
+\frac{(-1)^N(2N-1)!!}{2^N}I_N.
\]
All $I_m$ are positive, and the recurrence implies
$0<I_N<e^{-x^2}/(2x^{2N+1})$. The remainder in $\mathrm{erfc}$
has sign $(-1)^N$ and magnitude less than its first omitted term:
\[
|R_N|<\frac{e^{-x^2}}{\sqrt\pi x}
\frac{(2N-1)!!}{2^Nx^{2N}}.
\]
Finally $\mathrm{erf}=1-\mathrm{erfc}$ gives
\[
\boxed{\mathrm{erf}(x)\sim1-\frac{e^{-x^2}}{\sqrt\pi x}
\left[1-\frac1{2x^2}+\frac3{4x^4}-\frac{15}{8x^6}+\cdots\right]}.
\]
The term with index $k=n$ is
$(-1)^n(2n-1)!!/(2^nx^{2n})$, as requested.
The remainder bound proves the expansion for $x\to+\infty$.
For $x\to-\infty$ use $\mathrm{erf}(-x)=-\mathrm{erf}(x)$.

\Needspace{5\baselineskip}
\par\bigskip
\noindent\textbf{12.6.4}\quad The asymptotic expressions for the various
Bessel functions, Section 14.6, contain the series
\begin{align*}
P_{\nu}(z)&\sim1+\sum_{n=1}^{\infty}(-1)^{n}
\frac{\prod_{s=1}^{2n}[4\nu^{2}-(2s-1)^{2}]}{(2n)!(8z)^{2n}},\\
Q_{\nu}(z)&\sim\sum_{n=1}^{\infty}(-1)^{n+1}
\frac{\prod_{s=1}^{2n-1}[4\nu^{2}-(2s-1)^{2}]}{(2n-1)!(8z)^{2n-1}}.
\end{align*}
Show that these are asymptotic expansions \erratum{for fixed $\nu$ as
$z=x\to+\infty$; distinguish terminating half-integer orders from
divergent cases. Then specify a complex sector of validity.}
\solution
\erratum{Divergence at fixed $z$ does not by itself prove that a series
is asymptotic to a particular function. Also, for half-integer $\nu$ the
products terminate, so the blanket divergence argument in the manual
has exceptions. We first prove the asymptotic statement for fixed real
$\nu$ and $z=x\to+\infty$, then give its complex extension on a specified
sector.}
Put
\[
\mu=\nu^2-\frac14,\qquad
A_k=\frac{\prod_{s=1}^k[4\nu^2-(2s-1)^2]}{k!8^k},\quad A_0=1.
\]
The outgoing Bessel solution in Section 14.6 is written
\[
H_\nu^{(1)}(x)=\sqrt{\frac2{\pi x}}
 e^{i(x-\nu\pi/2-\pi/4)}w(x),\qquad w=P_\nu+iQ_\nu.
\]
Substituting $H=x^{-1/2}e^{ix}$ times a constant times $w$ into
$x^2H''+xH'+(x^2-\nu^2)H=0$ gives
\[
w''+2iw'-\frac\mu{x^2}w=0,\qquad w(x)\to1.
\]
Only the leading outgoing normalization $w\to1$ from Section 14.6
is used here; the higher coefficients and their remainder will be proved.
For $w\sim\sum_{k\ge0}c_kx^{-k}$, the coefficient of $x^{-k-2}$ is
\[
[k(k+1)-\mu]c_k-2i(k+1)c_{k+1}=0.
\]
Hence
\[
c_{k+1}=\frac{i[4\nu^2-(2k+1)^2]}{8(k+1)}c_k,\qquad
c_k=i^kA_k.
\]
The real even powers and imaginary odd powers are exactly the two
series in the question, including their signs.

Here is a remainder argument. For an integrable forcing $f$, define
\[
(Kf)(x)=\frac1{2i}\int_x^\infty[e^{2i(t-x)}-1]f(t)\,dt.
\]
Differentiating under the integral (the boundary kernel is zero) gives
\[
(Kf)'=-\int_x^\infty e^{2i(t-x)}f(t)\,dt,\qquad
(Kf)''=f(x)+2i\int_x^\infty e^{2i(t-x)}f(t)\,dt.
\]
Thus $(Kf)''+2i(Kf)'=f$, and $Kf\to0$ at infinity.
Apply this to $f=\mu w/x^2$, which is integrable because $w\to1$.
The difference $w-1-Kf$ solves $d''+2id'=0$, so $d=A+Be^{-2ix}$.
Since $d\to0$ on the positive axis, $A=B=0$.
This proves the Volterra equation
\[
w(x)=1+\frac\mu{2i}\int_x^\infty
 [e^{2i(t-x)}-1]\frac{w(t)}{t^2}dt.
\]
Its kernel has absolute value at most $1$ after division by $2i$.
On $x\ge R>|\mu|$, its operator $T$ has sup-norm at most
$|\mu|/R<1$, so the bounded outgoing solution is uniquely obtained by
successive substitution. Let $S_N=\sum_{k=0}^{N-1}c_kx^{-k}$.
The recurrence cancels every lower power, leaving precisely
\[
r_N:=S_N''+2iS_N'-\mu S_N/x^2
=[N(N-1)-\mu]c_{N-1}x^{-N-1}.
\]
Then $E_N=w-S_N$ satisfies
\[
E_N=T[E_N]-\frac1{2i}\int_x^\infty[e^{2i(t-x)}-1]r_N(t)dt.
\]
The explicit last integral is $O(x^{-N})$. In the weighted norm
$\|E\|_N=\sup_{x\ge R}x^N|E(x)|$,
\[
x^N|T[E](x)|
\le |\mu|\|E\|_Nx^N\int_x^\infty t^{-N-2}dt
\le\frac{|\mu|}{(N+1)R}\|E\|_N.
\]
The same contraction argument constructs $E_N=O(x^{-N})$; uniqueness
identifies $S_N+E_N$ with $w$. This is the required asymptotic remainder
estimate for every fixed truncation. Taking real and imaginary parts
proves the $P_\nu,Q_\nu$ expansions; retaining one extra term when
necessary gives the next even or odd inverse-power remainder.

For nonterminating cases the ratio of two consecutive $P$ terms is
\[
\left|\frac{[4\nu^2-(4n+1)^2][4\nu^2-(4n+3)^2]}
{(2n+1)(2n+2)(8x)^2}\right|\sim\frac{n^2}{x^2},
\]
and the $Q$ ratio has the same growth. Thus they diverge for fixed $x$,
while their finite truncations have the proved large-$x$ meaning.
If $\nu=\pm(j+1/2)$, $j=0,1,\ldots$, a numerator factor vanishes at $s=j+1$ and
all later coefficients are zero; the resulting finite expansions remain
valid, and in fact solve the transformed equation exactly.
\erratum{These terminating cases satisfy the usual fixed-order remainder
definition. They are excluded only if one imposes the textbook's extra
divergence convention, Eq. (12.98); see Exercise 12.6.5.}

The proof also supplies a complex sector, rather than assuming that an
expansion on the real axis automatically extends to all directions.
For $\Re z\ge0$, integrate along $t=z+u$, $u\ge0$. Then
\[
|e^{2iu}-1|\le2,\qquad
|z+u|\ge (|z|+u)/\sqrt2.
\]
Thus the unweighted operator norm on $|z|\ge R$ is at most
$2|\mu|/R$. In the weighted norm its bound is
\[
\frac{|\mu|\,2^{(N+2)/2}}{(N+1)R},
\]
and the residual integral is bounded by
$|[N(N-1)-\mu]c_{N-1}|\,2^{(N+1)/2}|z|^{-N}/N$.
Taking $R$ sufficiently large gives the same contraction and
$O(|z|^{-N})$ remainder, uniformly in this closed right half-plane.
The incoming solution obeys $w_-''-2iw_-'-\mu w_-/z^2=0$; its kernel
uses $e^{-2iu}-1$, with the same bound, and its coefficients are
$(-i)^kA_k$. Therefore $P=(w+w_-)/2$ selects the even terms and
$Q=(w-w_-)/(2i)$ selects the odd terms. These analytic solutions agree
with the Hankel normalizations on the positive axis and hence throughout
the sector. This argument also works for each fixed complex $\nu$,
since only $|\mu|$ enters the bounds. Wider continuation sectors require
the corresponding Hankel branches; no expansion across an unspecified
branch or Stokes boundary is being asserted.

\Needspace{5\baselineskip}
\par\bigskip
\noindent\textbf{12.6.5}\quad For $x>1$,
\[
\frac{1}{1+x}=\sum_{n=0}^{\infty}(-1)^{n}\frac{1}{x^{n+1}}.
\]
Test this series to see if it is an asymptotic series.
\solution
The answer depends on the stated convention. With the textbook's partial
sum through $x^{-m}$, $m\ge1$, finite geometric summation gives
\[
s_m(x)=\sum_{k=1}^m(-1)^{k-1}x^{-k}
=\frac{1-(-1/x)^m}{1+x},\qquad
R_m(x)=\frac{(-1)^m}{x^m(1+x)}.
\]
Consequently
\[
x^mR_m(x)=\frac{(-1)^m}{1+x}.
\]
For each fixed $m$ this tends to zero as $x\to+\infty$, satisfying
Eq. (12.97). For each fixed $x>1$, however, its absolute value is the
constant $1/(1+x)$; it does not diverge as $m\to\infty$.
Thus the extra requirement in Eq. (12.98) is not met:
\[
\boxed{\text{not an asymptotic series under the textbook's convention}.}
\]
This agrees with the manual. Its convergence follows independently from
the geometric ratio $|-1/x|<1$.

\erratum{Terminology qualification: the usual Poincar\'e definition uses
the fixed-order condition (12.97) without the additional divergence
requirement (12.98). The book itself acknowledges this alternative in
footnote 5 on p.583. Under that definition the same series is both
convergent and asymptotic; the manual's answer is correct under the
book's explicitly stricter convention.}
In that usual notation,
\[
\frac1{1+x}\sim\frac1x-\frac1{x^2}+\frac1{x^3}-\cdots,\qquad x\to+\infty.
\]
The exact remainder above proves the claim under either convention
without inferring asymptotic behavior merely from convergence or divergence.

\Needspace{5\baselineskip}
\par\bigskip
\noindent\textbf{12.6.6}\quad Derive the following Bernoulli-number asymptotic
series for the Euler-Mascheroni constant, defined in Eq.~(1.13):
\[
\gamma\sim\sum_{s=1}^{n}s^{-1}-\ln n-\frac{1}{2n}
+\sum_{k=1}^{\infty}\frac{B_{2k}}{(2k)n^{2k}}.
\]
Here $n$ plays the role of $x$.

\emph{Hint.}~Apply the Euler-Maclaurin integration formula to $f(x)=x^{-1}$
over the interval $[1,n]$ for $N=1,2,\ldots$.
\solution
Write $H_n=\sum_{s=1}^n1/s$. Applying Euler--Maclaurin to the finite
tail $\sum_{s=n+1}^M1/s$, with
$f^{(2k-1)}(x)=-(2k-1)!x^{-2k}$, gives
\[
H_M-H_n=\ln(M/n)+\frac1{2M}-\frac1{2n}
+\sum_{k=1}^N\frac{B_{2k}}{2k}
 (n^{-2k}-M^{-2k})+\text{remainder}.
\]
Let $M\to\infty$ and use the definition of $\gamma$. This produces
\[
\gamma=H_n-\ln n-\frac1{2n}
+\sum_{k=1}^N\frac{B_{2k}}{2k\,n^{2k}}+R_{n,N}.
\]
To prove the asymptotic nature, rather than merely write formal
corrections, we now derive a signed bound for $R_{n,N}$.
Define
\[
K(n)=\int_0^\infty e^{-nt}
 \left[\frac1t-\frac1{e^t-1}\right]dt.
\]
This integral is finite: the bracket tends to $1/2$ at zero and is
$O(1/t)$ at infinity. Its consecutive difference is
\begin{align*}
K(n)-K(n+1)
&=\int_0^\infty\frac{e^{-nt}-e^{-(n+1)t}}t dt
 -\int_0^\infty e^{-(n+1)t}dt\\
&=\ln(1+1/n)-\frac1{n+1}.
\end{align*}
The first equality used $(1-e^{-t})/(e^t-1)=e^{-t}$;
the logarithmic integral follows by integrating
$\int_n^{n+1}e^{-at}da$ first in $t$. This is also the consecutive
difference of $H_n-\ln n-\gamma$. The bracket in $K(n)$ is bounded,
say by $C$, so $|K(n)|\le C\int_0^\infty e^{-nt}dt=C/n\to0$.
Both sequences therefore tend to zero as $n\to\infty$,
so they are equal. Thus, with
$q(t)=1/(e^t-1)-1/t+1/2$,
\[
\gamma=H_n-\ln n-\frac1{2n}+\int_0^\infty e^{-nt}q(t)dt.
\]
The cotangent partial fractions, evaluated at $it/2$, give
\[
q(t)=\frac12\coth(t/2)-\frac1t
=2t\sum_{j=1}^{\infty}\frac1{(2\pi j)^2+t^2}.
\]
Expand each denominator by a \emph{finite} geometric identity through
$t^{2N-2}$. With $a_j=2\pi j$, this yields
\[
q(t)=\sum_{k=1}^N\frac{B_{2k}}{(2k)!}t^{2k-1}
+(-1)^N2t^{2N+1}\sum_{j=1}^{\infty}
\frac1{a_j^{2N}(a_j^2+t^2)}.
\]
The coefficients follow from the zeta/Bernoulli identity, and the last
sum is positive. Since $\int_0^\infty e^{-nt}t^m dt=m!/n^{m+1}$,
integration gives the stated correction terms and the rigorous bound
\[
\boxed{0<(-1)^N R_{n,N}
<\frac{|B_{2N+2}|}{(2N+2)n^{2N+2}}}.
\]
Indeed, replace $(a_j^2+t^2)^{-1}$ by $a_j^{-2}$ in the positive
remainder and use
$2(2N+1)!\zeta(2N+2)/(2\pi)^{2N+2}=|B_{2N+2}|/(2N+2)$.
For each fixed $N$ this remainder is smaller by two powers of $n$
than the last correction, proving the requested asymptotic series.
The coefficients grow factorially, so taking infinitely many terms at
fixed $n$ is not justified.

\Needspace{5\baselineskip}
\par\bigskip
\noindent\textbf{12.6.7}\quad Develop an asymptotic series for
\[
\int_{0}^{\infty}\frac{e^{-xv}}{\erratum{1+v^{2}}}\,dv.
\]
Take $x$ to be real and positive.
\par\nobreak\hfill\emph{ANS.}\enspace
$\dfrac{1}{x}-\dfrac{2!}{x^{3}}+\dfrac{4!}{x^{5}}-\cdots
+\dfrac{(-1)^{n}(2n)!}{x^{2n+1}}$.
\solution
\erratum{The published errata (Instructor's Manual, 6 April 2012,
textbook p.~585) replace the original denominator $(1+v^2)^2$ by
$1+v^2$. The corrected integral in the question now matches its printed
answer. We give both versions below to make the change explicit.}
For the corrected integral $I(x)$, use the exact finite identity
\[
\frac1{1+v^2}=\sum_{k=0}^N(-1)^kv^{2k}
+\frac{(-1)^{N+1}v^{2N+2}}{1+v^2}.
\]
Multiply by $e^{-xv}$ and integrate for $x>0$. With $u=xv$,
\[
\int_0^\infty e^{-xv}v^{2k}dv
=x^{-2k-1}\int_0^\infty e^{-u}u^{2k}du
=\frac{(2k)!}{x^{2k+1}},
\]
where repeated integration by parts evaluates the last integral.
Consequently
\[
\boxed{I(x)\sim\sum_{k=0}^{\infty}\frac{(-1)^k(2k)!}{x^{2k+1}}
=\frac1x-\frac{2!}{x^3}+\frac{4!}{x^5}-\cdots}.
\]
The actual remainder after $k=N$ is
\[
R_N=(-1)^{N+1}\int_0^\infty
 \frac{e^{-xv}v^{2N+2}}{1+v^2}dv,\qquad
0<(-1)^{N+1}R_N<\frac{(2N+2)!}{x^{2N+3}}.
\]
This proves the asymptotic statement without interchanging an infinite
geometric series outside its disk of convergence.

For comparison, the \emph{original} squared-denominator integral has
\[
\int_0^\infty\frac{e^{-xv}}{(1+v^2)^2}dv
\sim\boxed{\sum_{k=0}^{\infty}
\frac{(-1)^k(k+1)(2k)!}{x^{2k+1}}}
=\frac1x-\frac{4}{x^3}+\frac{72}{x^5}-\cdots.
\]
Its coefficients follow from $(1+w)^{-2}=\sum(-1)^k(k+1)w^k$.
More precisely, the finite remainder is
$(-1)^{N+1}w^{N+1}[(N+2)+(N+1)w]/(1+w)^2$.
For $w=v^2\ge0$ its magnitude is at most $(N+2)v^{2N+2}$,
so integration bounds the original integral's remainder by
$(N+2)(2N+2)!/x^{2N+3}$. The factor $k+1$ shows exactly why the
original problem and printed answer were inconsistent.

\Needspace{8\baselineskip}
\subsection*{12.7 Method of Steepest Descents}

We present here a rather small number of exercises on the method of steepest
descents. Several additional exercises appear elsewhere in this book, in
particular in Section 14.6, where the technique is applied to the contour
integral representations of Bessel functions.

\Needspace{5\baselineskip}
\par\bigskip
\noindent\textbf{12.7.1}\quad \erratum{Let $F$ be nonconstant and analytic on a connected
open region. Derive the correct circle mean of $|F|^2$ about $z_0$
from the Taylor coefficients of $F$. Prove that $|F|$ has no interior
local maximum and no positive interior local minimum. In particular,
if $F$ never vanishes, it has no interior modulus extremum.}
\solution
\erratum{Both the asserted mean-value identity for $|F|^2$ and the
unqualified absence of minima are false. For $F(z)=z$ at $z_0=0$,
the circle average is $r^2$, whereas $|F(0)|^2=0$ is a minimum.
The nonconstant, nonzero case used in the preceding text needs a
corrected argument.}
Expand in a disk about $z_0$ lying in the analytic region:
\[
F(z_0+re^{i\theta})=\sum_{n=0}^{\infty}a_nr^ne^{in\theta}.
\]
Uniform convergence on a smaller circle permits multiplication and
integration term by term. Since
$(2\pi)^{-1}\int_0^{2\pi}e^{i(n-m)\theta}d\theta=\delta_{nm}$,
\begin{align*}
\frac1{2\pi}\int_0^{2\pi}|F(z_0+re^{i\theta})|^2d\theta
&=\sum_{n,m\ge0}a_n\overline{a_m}r^{n+m}\delta_{nm}\\
&=\boxed{|F(z_0)|^2+\sum_{n=1}^{\infty}|a_n|^2r^{2n}}.
\end{align*}
At least one $a_n$, $n\ge1$, is nonzero: otherwise the Taylor series
would make $F$ constant near $z_0$, and the identity theorem would make
it constant throughout the connected region. Thus the mean strictly exceeds the center value on every
sufficiently small circle. A local maximum would make every value on
that circle no larger than the center, a contradiction. Therefore
$|F|$ has no interior local maximum unless $F$ is constant.

If $F(z_0)\ne0$, continuity gives a smaller disk where $F$ never
vanishes. Apply the same maximum argument to the analytic function
$1/F$. A local minimum of $|F|$ would be a local maximum of $|1/F|$,
again impossible for nonconstant $F$. Thus a nonconstant analytic
function has no positive interior local minimum; zeros are allowed
minima. A constant function is the usual exception to both statements.

For comparison, the true mean-value identities are
$\langle F\rangle=F(z_0)$, and, on a zero-free disk,
$\langle\log|F|\rangle=\log|F(z_0)|$. The latter follows by taking
the real part of the mean-value formula for a local analytic logarithm.
It directly proves the no-extremum assertion in the zero-free setting
relevant to the saddle-point discussion.

\Needspace{5\baselineskip}
\par\bigskip
\noindent\textbf{12.7.2}\quad Find the steepest path and leading asymptotic
expansion \erratum{as $s\to+\infty$} for the Fresnel integrals
$\displaystyle\int_{0}^{s}\cos x^{2}\,dx$,\quad
$\displaystyle\int_{0}^{s}\sin x^{2}\,dx$.

\emph{Hint.}~Use $\displaystyle\int_{0}^{1}e^{itz^{2}}\,dz$.
\solution
Put $t=s^2>0$ and scale $x=sz$:
\[
\int_0^s e^{ix^2}dx=sJ(t),\qquad J(t)=\int_0^1e^{itz^2}dz.
\]
For $f(z)=iz^2$, the saddle is $z=0$, with $f''(0)=2i$.
Writing $z=re^{i\theta}$ gives
$\Re(it z^2)=-tr^2\sin2\theta$ and
$\Im(it z^2)=tr^2\cos2\theta$. The descending ray from zero is
therefore $\theta=\pi/4$; on it the exponential becomes $e^{-tr^2}$.
The endpoint $z=1$ has a descending constant-phase path
\[
z=\sqrt{1+iv},\quad 0\le v<\infty,\qquad
it z^2=it-tv,\quad dz=\frac{i\,dv}{2\sqrt{1+iv}}.
\]
Both paths approach the same direction at infinity. Deforming the finite
interval to their difference is legitimate because the integrand is
entire and the joining segment at infinity has exponentially vanishing
integral. Hence the exact contour decomposition is
\[
J(t)=e^{i\pi/4}\int_0^\infty e^{-tr^2}dr
-\frac i2e^{it}\int_0^\infty\frac{e^{-tv}}{\sqrt{1+iv}}dv.
\]
The first integral is the endpoint half-Gaussian
$\sqrt\pi/(2\sqrt t)$. For the second,
$(1+iv)^{-1/2}=1-iv/2-3v^2/8+O(v^3)$ near zero.
Termwise integration of any finite Taylor truncation gives
\[
J(t)=\frac{e^{i\pi/4}\sqrt\pi}{2\sqrt t}
-\frac{ie^{it}}{2t}-\frac{e^{it}}{4t^2}+O(t^{-3}).
\]
To justify the error, split the Laplace integral at a fixed small
$v=\delta$: its Taylor remainder contributes $O(\int_0^\infty
v^2e^{-tv}dv)=O(t^{-3})$ at the order displayed, and its remote tail
is exponentially small.
Multiplying by $s$ and taking real and imaginary parts yields
\begin{align*}
\boxed{\int_0^s\cos x^2dx}
&\boxed{=\sqrt{\frac\pi8}+\frac{\sin s^2}{2s}
-\frac{\cos s^2}{4s^3}+O(s^{-5}),}\\
\boxed{\int_0^s\sin x^2dx}
&\boxed{=\sqrt{\frac\pi8}-\frac{\cos s^2}{2s}
-\frac{\sin s^2}{4s^3}+O(s^{-5}).}
\end{align*}
\erratum{The saddle is an endpoint, so it contributes half of a full
Gaussian integral. The manual's limiting value is twice the correct
$\sqrt{\pi/8}$.}

\Needspace{5\baselineskip}
\par\bigskip
\noindent\textbf{12.7.3}\quad Show that the formula
\[
\Gamma(1+s)\approx\sqrt{2\pi s}\,s^{s}e^{-s}
\]
holds for complex values of $s$ (with $\Re(s)$ large and positive).

\emph{Hint.}~This involves assigning a phase to $s$ and then demanding that
$\Im[sf(z)]$ be constant in the vicinity of the saddle point.
\solution
Use the principal $\operatorname{Log}s$ for $\Re s>0$.
Starting with $\Gamma(1+s)=\int_0^\infty e^{-u}u^sdu$, set $u=sz$.
The resulting ray can be rotated to the positive $z$ axis: the rotation
stays inside $|\arg u|<\pi/2$, the exponential decays at infinity,
and $u^s$ is integrable at zero since $\Re s>0$.
Along the deformed rays $\operatorname{Log}u=\operatorname{Log}s+
\operatorname{Log}z$, so $u^s=s^sz^s$ with consistent branches. Thus
\[
\Gamma(1+s)=s^{s+1}\int_0^\infty e^{s f(z)}dz,\qquad
f(z)=\operatorname{Log}z-z.
\]
The saddle is $z_0=1$, because $f'=1/z-1$; there
$f(1)=-1$ and $f''(1)=-1$. Write $s=\rho e^{i\phi}$.
Near the saddle,
\[
sf(z)=-s-\frac s2(z-1)^2+\frac s3(z-1)^3+\cdots.
\]
To make the quadratic decrement real and negative, choose
\[
z-1=e^{-i\phi/2}\frac{v}{\sqrt\rho}=\frac v{\sqrt s},\qquad v\in\mathbb R.
\]
This is the local steepest path and keeps $\Im[sf(z)]$ constant to
quadratic order. Since $dz=dv/\sqrt s$, its leading Gaussian
contribution is
\[
\int e^{sf(z)}dz
\sim\frac{e^{-s}}{\sqrt s}\int_{-\infty}^{\infty}e^{-v^2/2}dv
=\frac{\sqrt{2\pi}e^{-s}}{\sqrt s}.
\]
Therefore, with the same branch for all powers,
\[
\boxed{\Gamma(1+s)\sim\sqrt{2\pi s}\,s^s e^{-s}}.
\]
The factor $e^{-i\phi/2}$ from the path is essential; replacing the
Gaussian width by $1/\sqrt{|s|}$ would lose that phase.

For completeness, the approximation remains valid throughout the
stated limit $\Re s\to+\infty$, even if the phase changes with $s$.
A useful uniform remainder follows from Euler--Maclaurin. The Euler
limit for $\Gamma(s)$ gives
\[
\log\Gamma(s)=\lim_{M\to\infty}
\left[\log M!+s\log M-\sum_{k=0}^M\operatorname{Log}(s+k)\right].
\]
Apply Euler--Maclaurin to $h(t)=\operatorname{Log}(s+t)$:
\[
\sum_{k=0}^Mh(k)=\int_0^Mh(t)dt+\frac{h(0)+h(M)}2
+\frac{h'(M)-h'(0)}{12}
-\frac12\int_0^M\widetilde B_2(t)h''(t)dt,
\]
where $\widetilde B_2(t)=B_2(\{t\})$ is periodic. Using
$h'=1/(s+t)$, $h''=-1/(s+t)^2$, and
$\int h=(s+t)\operatorname{Log}(s+t)-(s+t)$, the finite sum is
\begin{align*}
\sum_{k=0}^M\operatorname{Log}(s+k)
={}&(s+M+\tfrac12)\operatorname{Log}(s+M)-M\\
&-(s-\tfrac12)\operatorname{Log}s
 +\frac1{12(s+M)}-\frac1{12s}\\
&+\frac12\int_0^M\frac{\widetilde B_2(t)}{(s+t)^2}dt.
\end{align*}
The real Gaussian calculation gives
$\log M!=(M+\tfrac12)\log M-M+\tfrac12\log(2\pi)+o(1)$.
After inserting both formulas in the Euler limit, the remaining
$M$-dependent logarithms are
\[
(M+s+\tfrac12)\bigl[\log M-\operatorname{Log}(M+s)\bigr]
=-(M+s+\tfrac12)\operatorname{Log}(1+s/M)\longrightarrow-s
\]
for each fixed $s$ with $\Re s>0$; also $1/(M+s)\to0$.
Adding $\operatorname{Log}s$ to pass from $\Gamma(s)$ to $\Gamma(1+s)$
therefore yields
\[
\log\Gamma(1+s)=(s+\tfrac12)\operatorname{Log}s-s
+\tfrac12\ln2\pi+R(s),\qquad
R(s)=\frac1{12s}-\frac12\int_0^\infty
\frac{\widetilde B_2(t)}{(s+t)^2}dt.
\]
Since $|\widetilde B_2|\le1/6$ and $|s+t|\ge\Re s+t$,
\[
|R(s)|\le\frac1{12|s|}+\frac1{12\Re s}
\le\frac1{6\Re s}\longrightarrow0.
\]
Exponentiation proves that the ratio to the boxed expression tends to
one for the entire requested right-half-plane limit.

\Needspace{8\baselineskip}
\subsection*{12.8 Dispersion Relations}

\Needspace{5\baselineskip}
\par\bigskip
\noindent\textbf{12.8.1}\quad Assume that the function $f(z)$ satisfies the
conditions for the dispersion relations. In addition, assume that
$f(z)=f^{*}(z^{*})$, i.e., that it meets the conditions of the Schwarz
reflection principle, Eq.~(11.127). Show that $f(z)$ is identically zero.
\solution
The dispersion assumptions are analyticity in the upper half-plane and
across the real axis, with sufficient decay for the large arc to vanish.
On the real axis the additional reflection condition gives
$f(x)=f(x)^*$, so $v(x)=\Im f(x)=0$.
The first Hilbert-transform equation is
\[
u(x_0)=\frac1\pi\operatorname{PV}\int_{-\infty}^{\infty}
\frac{v(x)}{x-x_0}dx.
\]
Its integrand vanishes identically, hence $u(x_0)=0$ for every real
$x_0$ too. Thus $f=0$ on the real axis. Because that axis lies in a
region through which $f$ is analytic, the identity theorem implies
\[
\boxed{f(z)\equiv0}
\]
on the connected analytic upper-half-plane domain. Equivalently,
Schwarz reflection would extend $f$ to an entire function with decay
in both half-planes, and Liouville's theorem gives the same result.

\Needspace{5\baselineskip}
\par\bigskip
\noindent\textbf{12.8.2}\quad For $f(z)$ such that we may replace the closed
contour of the Cauchy integral formula by an integral over the real axis we
have
\[
f(x_{0})=\frac{1}{2\pi i}\left\{
\int_{-\infty}^{x_{0}-\delta}\frac{f(x)}{x-x_{0}}\,dx
+\int_{x_{0}+\delta}^{\infty}\frac{f(x)}{x-x_{0}}\,dx\right\}
+\frac{1}{2\pi i}\int_{C}\frac{f(x)}{x-x_{0}}\,dx.
\]
Here we take $C$ to be a small semicircle about $x_{0}$ in the lower
half-plane. Show that \erratum{as $\delta\downarrow0$} the formula for $f(x_{0})$ reduces to
\[
f(x_{0})=\frac{1}{\pi i}\erratum{\operatorname{PV}}\int_{-\infty}^{\infty}\frac{f(x)}{x-x_{0}}\,dx,
\]
which is Eq.~(12.117).
\solution
Parameterize the lower semicircle from its left endpoint to its right
endpoint as
\[
z=x_0+\delta e^{i\theta},\quad -\pi\le\theta\le0,\qquad
dz=i\delta e^{i\theta}d\theta.
\]
Its orientation is counterclockwise around the lower half of the small
circle, so
\[
\int_C\frac{f(z)}{z-x_0}dz
=i\int_{-\pi}^0f(x_0+\delta e^{i\theta})d\theta
\longrightarrow i\pi f(x_0).
\]
The two real integrals together tend to their Cauchy principal value.
Taking $\delta\to0$ in the formula in the question therefore gives
\[
f(x_0)=\frac1{2\pi i}\operatorname{PV}\int_{-\infty}^{\infty}
\frac{f(x)}{x-x_0}dx+\frac12f(x_0).
\]
Subtract the half contribution and multiply by two:
\[
\boxed{f(x_0)=\frac1{\pi i}\operatorname{PV}
\int_{-\infty}^{\infty}\frac{f(x)}{x-x_0}dx}.
\]
\erratum{The real-axis integral in the requested formula must be read as
a Cauchy principal value; an ordinary improper integral at the pole
need not exist. The source uses a cut integral sign for this convention.}

\Needspace{5\baselineskip}
\par\bigskip
\noindent\textbf{12.8.3}\quad
\begin{enumerate}
\item[(a)] The function $f(z)=e^{iz}$ does not vanish at the endpoints of the
  range of $\arg z$, \erratum{$0$} and $\pi$. Show, with the help of Jordan's lemma,
  Eq.~(11.102), that Eq.~(12.116) still holds.
\item[(b)] For $f(z)=e^{iz}$ verify by direct integration the dispersion
  relations, Eq.~(12.117) or Eqs.~(12.118).
\end{enumerate}
\solution
\erratum{The printed lower endpoint \texttt{a} is $0$: the contour is the
upper semicircle $0\le\arg z\le\pi$.}
(a) For $z_0$ fixed in the upper half-plane, put $z=Re^{i\theta}$ on
the large semicircle, $0\le\theta\le\pi$. Then
\[
\left|\int_{\rm arc}\frac{e^{iz}}{z-z_0}dz\right|
\le\frac R{R-|z_0|}\int_0^\pi e^{-R\sin\theta}d\theta.
\]
On $0\le\theta\le\pi/2$, $\sin\theta\ge2\theta/\pi$; symmetry
therefore bounds the last integral by
$2\int_0^{\pi/2}e^{-2R\theta/\pi}d\theta\le\pi/R$.
The arc contribution is at most $\pi/(R-|z_0|)\to0$, despite the
nonvanishing endpoint values. Cauchy's formula consequently gives
\[
\boxed{e^{iz_0}=\frac1{2\pi i}\int_{-\infty}^{\infty}
\frac{e^{ix}}{x-z_0}dx},\qquad\Im z_0>0.
\]

(b) Set $t=x-x_0$. The symmetric principal value is
\[
\operatorname{PV}\int_{-\infty}^{\infty}\frac{e^{ix}}{x-x_0}dx
=e^{ix_0}\left[\operatorname{PV}\int_{-\infty}^{\infty}
\frac{\cos t}{t}dt+i\int_{-\infty}^{\infty}\frac{\sin t}{t}dt\right].
\]
The cosine integral is zero by oddness. To evaluate the sine integral,
first insert $e^{-\epsilon t}$ on $t>0$. If
$A(b)=\int_0^\infty e^{-\epsilon t}\sin(bt)dt/t$, then
$A'(b)=\epsilon/(\epsilon^2+b^2)$ and $A(0)=0$, so
$A(1)=\arctan(1/\epsilon)\to\pi/2$.
The damping can be removed uniformly in the tail: integration by parts
against a bounded primitive of $\sin t$ gives
$|\int_A^\infty e^{-\epsilon t}\sin t\,dt/t|\le2/A$ for every
$\epsilon\ge0$. On $[0,A]$, dominated convergence applies.
Thus the whole-line sine integral is $\pi$, and
\[
\operatorname{PV}\int_{-\infty}^{\infty}\frac{e^{ix}}{x-x_0}dx
=i\pi e^{ix_0}.
\]
Dividing by $i\pi$ verifies Eq. (12.117). Separating components also gives
\[
\frac1\pi\operatorname{PV}\int\frac{\sin x}{x-x_0}dx=\cos x_0,\qquad
-\frac1\pi\operatorname{PV}\int\frac{\cos x}{x-x_0}dx=\sin x_0,
\]
with both integrals over the whole real line, verifying both equations
in (12.118).

\Needspace{5\baselineskip}
\par\bigskip
\noindent\textbf{12.8.4}\quad With $f(x)=u(x)+iv(x)$ and $f(x)=f^{*}(-x)$,
\erratum{assume the dispersion relations. For (a), require
$v\in C^1$ and $v^{(j)}(x)=O(x^{-2-\epsilon-j})$; for (b), require
$u\in C^1$ and $u^{(j)}(x)=O(x^{-1-\epsilon-j})$, $j=0,1$,
with $0<\epsilon<1$.} Show that as $x_0\to+\infty$,
\begin{enumerate}
\item[(a)] $u(x_{0})\erratum{=}-\dfrac{2}{\pi x_{0}^{2}}
  \displaystyle\int_{0}^{\infty}xv(x)\,dx\erratum{+o(x_0^{-2})}$,
\item[(b)] $v(x_{0})\erratum{=}\dfrac{2}{\pi x_{0}}
  \displaystyle\int_{0}^{\infty}u(x)\,dx\erratum{+o(x_0^{-1})}$.
\end{enumerate}
In quantum mechanics relations of this form are often called \emph{sum rules}.
\solution
The crossing condition implies $u(-x)=u(x)$ and $v(-x)=-v(x)$.
Split the Hilbert integrals into positive and negative halves, substitute
$x\mapsto-x$ on the latter, and combine denominators. With $a=x_0>0$,
\[
u(a)=\frac2\pi\operatorname{PV}\int_0^\infty
\frac{xv(x)}{x^2-a^2}dx,\qquad
v(a)=-\frac{2a}\pi\operatorname{PV}\int_0^\infty
\frac{u(x)}{x^2-a^2}dx.
\]
\erratum{The displayed moments and the passage to large $a$ require
additional decay/tail control. Crossing and the ordinary dispersion
assumptions alone do not suffice. Also, a zero leading moment gives a
little-$o$ conclusion, not an equivalence to zero.}
For example $f(z)=1/(1-iz)$ is analytic above the real axis, decays,
and satisfies crossing, but $v(x)=x/(1+x^2)$ makes
$\int_0^\infty xv(x)dx$ divergent.

Here are sufficient assumptions and a justified expansion. For (a),
assume $v,v'$ are continuous and, for some $0<\epsilon<1$,
$v^{(j)}(x)=O(x^{-2-\epsilon-j})$, $j=0,1$.
For (b), assume instead
$u^{(j)}(x)=O(x^{-1-\epsilon-j})$, $j=0,1$.
Each condition is imposed on the function in its corresponding integral.
The exact kernel decompositions give
\begin{align*}
u(a)&=-\frac2{\pi a^2}\int_0^\infty xv(x)dx
 +\frac2{\pi a^2}\operatorname{PV}\int_0^\infty
       \frac{x^3v(x)}{x^2-a^2}dx,\\
v(a)&=\frac2{\pi a}\int_0^\infty u(x)dx
 -\frac2{\pi a}\operatorname{PV}\int_0^\infty
       \frac{x^2u(x)}{x^2-a^2}dx.
\end{align*}
Both final integrals tend to zero. To see the control near the moving
pole, put $h=x^3v$ in (a) or $h=x^2u$ in (b), and split the range
at $a/2,3a/2$. On the middle interval write
$g(x)=h(x)/(x+a)$, so the principal value is
\[
\operatorname{PV}\int_{a/2}^{3a/2}\frac{g(x)}{x-a}dx
=\int_{a/2}^{3a/2}\frac{g(x)-g(a)}{x-a}dx.
\]
The constant term cancels by symmetry; $g'=O(a^{-1-\epsilon})$
bounds this integral by $O(a^{-\epsilon})$.
On the lower interval, $|x^2-a^2|\ge3a^2/4$ and
$h(x)=O(x^{1-\epsilon})$ at large $x$, so its contribution is
$O(a^{-\epsilon})+O(a^{-2})$. On the upper interval the denominator
is bounded below by a constant times $x^2$, giving
$O(\int_{3a/2}^\infty x^{-1-\epsilon}dx)=O(a^{-\epsilon})$.
It follows that
\[
\boxed{u(a)=-\frac2{\pi a^2}\int_0^\infty xv(x)dx+o(a^{-2})},
\]
\[
\boxed{v(a)=\frac2{\pi a}\int_0^\infty u(x)dx+o(a^{-1})}.
\]
When the respective moment is nonzero, these imply exactly the two
originally printed $\sim$ relations. When it vanishes, a further moment is
needed to find the first nonzero term.

\Needspace{5\baselineskip}
\par\bigskip
\noindent\textbf{12.8.5}\quad
\begin{enumerate}
\item[(a)] Given the integral equation (valid for all real $x_{0}$)
  \[
  \frac{1}{1+x_{0}^{2}}
  =\frac{1}{\pi}\erratum{\operatorname{PV}}\int_{-\infty}^{\infty}\frac{u(x)}{x-x_{0}}\,dx,
  \]
  use Hilbert transforms to determine $u(x_{0})$.
\item[(b)] Verify that the $u(x_{0})$ found as your answer to part (a)
  actually satisfies the integral equation.
\item[(c)] From $f(z)|_{y=0}=u(x)+iv(x)$, replace $x$ by $z$ and determine
  $f(z)$. Verify that the conditions for the Hilbert transforms are
  satisfied.
\item[(d)] Are the crossing conditions satisfied?
\end{enumerate}
\par\nobreak\hfill\emph{ANS.}\enspace
(a) $u(x_{0})=\dfrac{x_{0}}{1+x_{0}^{2}}$,\quad
(c) $f(z)=(z+i)^{-1}$.
\solution
\erratum{The principal-value symbol restores the cut integral sign
in the source; it is necessary at $x=x_0$.} All integrals below have
that meaning. Use the convention
\[
(Hg)(a)=\frac1\pi\operatorname{PV}\int_{-\infty}^{\infty}
\frac{g(x)}{x-a}dx,\qquad u=Hv,\quad v=-Hu.
\]
(a) The given equation is $Hu=1/(1+a^2)$, so
$v(a)=-1/(1+a^2)$ and $u=Hv$. To evaluate the latter transform, use
\[
\frac1{(1+x^2)(x-a)}
=\frac1{1+a^2}\left[\frac1{x-a}-\frac{x+a}{1+x^2}\right].
\]
The principal value of $1/(x-a)$ on the whole real line is zero;
$x/(1+x^2)$ contributes zero by symmetric odd integration; and
$\int dx/(1+x^2)=\pi$. Therefore
\[
u(a)=-\frac1\pi\operatorname{PV}\int
\frac{dx}{(1+x^2)(x-a)}=\boxed{\frac a{1+a^2}}.
\]
Within the decaying Hilbert-transform class, inversion $H^2=-1$
fixes this solution uniquely.

(b) Check the original transform directly with another decomposition:
\[
\frac{x}{(1+x^2)(x-a)}
=\frac1{1+a^2}\left[\frac a{x-a}+\frac{1-ax}{1+x^2}\right].
\]
Its integral is $\pi/(1+a^2)$ by the same three evaluations. Division
by $\pi$ reproduces the given equation, including its positive sign.

(c) On the real axis,
\[
f(x)=u+iv=\frac{x-i}{x^2+1}=\frac1{x+i}.
\]
Thus $\boxed{f(z)=1/(z+i)}$. Its only pole is at $-i$, below the
real axis, so it is analytic in and above the real axis. Also
$f(z)=O(1/z)$ there. For the Cauchy kernel $f(z)/(z-z_0)$ the large
arc is $O(1/R)$, verifying the needed dispersion condition.

(d) The standard crossing conditions are not satisfied: $u$ is odd,
whereas $v$ is even. In fact
\[
f(-x)=\frac1{-x+i}=-\frac1{x-i}=-f(x)^*,
\]
with a minus sign instead of the requested $f(-x)=f(x)^*$.

\Needspace{5\baselineskip}
\par\bigskip
\noindent\textbf{12.8.6}\quad
\erratum{Assume the full-frequency dispersion hypotheses for
$f(\omega)=n^2(\omega)-1$, including analyticity above the real axis
and high-frequency decay. In (a), constancy means on the entire real
frequency axis.}
\begin{enumerate}
\item[(a)] If the real part of the complex index of refraction (squared) is
  constant (no optical dispersion), show that the imaginary part is zero (no
  absorption).
\item[(b)] Conversely, if there is absorption, show that there must be
  dispersion. In other words, if the imaginary part of $n^{2}-1$ is not
  zero, show that the real part of $n^{2}-1$ is not constant.
\end{enumerate}
\solution
\erratum{The conclusion uses the dispersion hypotheses on
$f(\omega)=n^2(\omega)-1$, including its decay at high frequency,
and constancy over the entire frequency axis. Constancy over only a
finite frequency interval is not enough for this argument.}
Write $f=u+iv$. (a) If the real part of $n^2$ is constant, then
$u=\Re(n^2-1)=K$ is constant. Since $f(\omega)\to0$ at high frequency,
$K=0$. The second Hilbert-transform equation then gives
\[
v(\omega_0)=-\frac1\pi\operatorname{PV}\int_{-\infty}^{\infty}
\frac{u(\omega)}{\omega-\omega_0}d\omega=0.
\]
Thus $\boxed{\Im n^2=0}$, meaning no absorption under the stated model.
The constant real value of $n^2$ is in fact one because of the
high-frequency normalization.

(b) This is the contrapositive of (a). If $v$ is not identically zero
and $u$ were constant, the preceding argument would force $v=0$, a
contradiction. Hence
\[
\boxed{\Im(n^2-1)\not\equiv0\quad\Longrightarrow\quad
\Re(n^2-1)\text{ is not constant}}.
\]
This asserts the existence of dispersion somewhere in the full spectrum;
it does not demand that its size be appreciable at every frequency.

\Needspace{5\baselineskip}
\par\bigskip
\noindent\textbf{12.8.7}\quad Given $u(x)=x/(x^{2}+1)$ and
$v(x)=-1/(x^{2}+1)$, show by direct evaluation of each integral that
\[
\int_{-\infty}^{\infty}|u(x)|^{2}\,dx
=\int_{-\infty}^{\infty}|v(x)|^{2}\,dx.
\]
\par\nobreak\hfill\emph{ANS.}\enspace
$\displaystyle\int_{-\infty}^{\infty}|u(x)|^{2}\,dx
=\int_{-\infty}^{\infty}|v(x)|^{2}\,dx=\frac{\pi}{2}$.
\solution
Use $x=\tan\theta$, $dx=\sec^2\theta\,d\theta$, with
$-\pi/2<\theta<\pi/2$. For the real functions in the question,
\[
\int_{-\infty}^{\infty}|u|^2dx
=\int_{-\pi/2}^{\pi/2}
 \frac{\tan^2\theta}{\sec^4\theta}\sec^2\theta\,d\theta
=\int_{-\pi/2}^{\pi/2}\sin^2\theta\,d\theta,
\]
\[
\int_{-\infty}^{\infty}|v|^2dx
=\int_{-\pi/2}^{\pi/2}
 \frac1{\sec^4\theta}\sec^2\theta\,d\theta
=\int_{-\pi/2}^{\pi/2}\cos^2\theta\,d\theta.
\]
With $\sin^2\theta=(1-\cos2\theta)/2$ and
$\cos^2\theta=(1+\cos2\theta)/2$, the oscillatory endpoint term
$\sin2\theta/4$ vanishes at both limits. Each integral is therefore
\[
\boxed{\pi/2}.
\]
This directly checks the Parseval equality without assuming it.

\Needspace{5\baselineskip}
\par\bigskip
\noindent\textbf{12.8.8}\quad Take $u(x)=\delta(x)$, a delta function, and
assume that the Hilbert transform equations hold.
\erratum{Interpret both displayed kernel integrals as distributions:
replace each reciprocal $1/q$ by $q/(q^2+\epsilon^2)$, use the same
$\epsilon>0$ for both factors, integrate, and then take
$\epsilon\downarrow0$ against smooth compactly supported test functions.}
\begin{enumerate}
\item[(a)] Show that
  \[
  \delta(w)=\frac{1}{\pi^{2}}\int_{-\infty}^{\infty}\frac{dy}{y(y-w)}.
  \]
\item[(b)] With changes of variables $w=s-t$ and $x=s-y$, transform the
  $\delta$ representation of part (a) into
  \[
  \delta(s-t)=\frac{1}{\pi^{2}}
  \int_{-\infty}^{\infty}\frac{dx}{(x-s)(x-t)}.
  \]
\end{enumerate}
\emph{Note.}~The $\delta$ function is discussed in Section 1.11.
\solution
\erratum{These formulas are distributional kernel identities. They are
not ordinary improper integrals, nor does the pointwise principal value
at coincident poles exist. A consistent regularization is essential.}
With $u(x)=\delta(x)$, the second Hilbert relation gives
\[
v(y)=-\frac1\pi\operatorname{PV}\int\frac{\delta(x)}{x-y}dx
=\frac1\pi\operatorname{PV}\frac1y.
\]
The first relation $u=Hv$ therefore formally gives the expression in
(a). To specify and prove its meaning, regularize both kernels by the
same positive $\epsilon$ and define
\[
D_\epsilon(w)=\frac1{\pi^2}\int_{-\infty}^{\infty}
\frac{y(y-w)}{(y^2+\epsilon^2)((y-w)^2+\epsilon^2)}dy.
\]
For real $w\ne0$, close the rational integrand in the upper half-plane.
The arc vanishes as $O(1/R)$, and the two residues before division by
$\pi^2$ are
\[
\frac{-w+i\epsilon}{2w(w-2i\epsilon)},\qquad
\frac{w+i\epsilon}{2w(w+2i\epsilon)}.
\]
Their sum is $-i\epsilon/(w^2+4\epsilon^2)$, so
\[
\boxed{D_\epsilon(w)=\frac{2\epsilon}{\pi(w^2+4\epsilon^2)}}.
\]
The formula extends to $w=0$ by continuity; directly the integral there
is $1/(2\pi\epsilon)$, the same value. For every smooth compactly
supported test function $\varphi$, substitute $w=2\epsilon t$:
\[
\int_{-\infty}^{\infty}D_\epsilon(w)\varphi(w)dw
=\int_{-\infty}^{\infty}\frac{\varphi(2\epsilon t)}{\pi(1+t^2)}dt
\longrightarrow\varphi(0).
\]
Dominated convergence applies because $\varphi$ is bounded and the
Cauchy density has integral one. Therefore $D_\epsilon\to\delta$ in
the distributional sense, establishing the precise meaning of
\[
\boxed{\delta(w)=\frac1{\pi^2}\int\frac{dy}{y(y-w)}}.
\]
For $w\ne0$ the separately interpreted principal values would give
zero, while at $w=0$ the integral of $1/y^2$ diverges; this is why a
pointwise calculation cannot replace the distributional limit.

(b) Set $w=s-t$ and $x=s-y$. Then $y=s-x$, $y-w=t-x$, and
$dy=-dx$. Reversing the limits removes the minus sign, while
\[
y(y-w)=(s-x)(t-x)=(x-s)(x-t).
\]
Thus
\[
\boxed{\delta(s-t)=\frac1{\pi^2}
\int_{-\infty}^{\infty}\frac{dx}{(x-s)(x-t)}}
\]
with the same regularization: its finite-$\epsilon$ kernel is
$2\epsilon/[\pi((s-t)^2+4\epsilon^2)]$. This also supplies the
consistent kernel identity used in the Parseval calculation.
